% =====================================================================
% Chapitre 2 -- Le flux des formalisations
% Sources : FACTS.md ; papers/ar_dipolar_exciton, ar_yukawa_crystal, ar_keldysh_tmd,
% prog_cm_bernstein_library, prog_sandier_serfaty_audit, p2_sandier_serfaty_lean, p2_3d_feasibility ;
% Lean : formal_cert/Autoresearch/{DipolarExciton*,Yukawa*,KeldyshCM*}.lean,
% formal_cert/Programmes/CompletelyMonotone/, formal_cert/Programmes/SandierSerfaty/SignedEndpoints.lean.
% Les extraits Lean sont VERBATIM ; le numéro de ligne renvoie au fichier source.
% =====================================================================

% Style de listings local au chapitre : la liste « literate » de main.tex, recopiée telle quelle,
% complétée par quelques symboles présents dans les extraits de ce chapitre (β, ∫, ⁻, ᵐ, ι, ᵢ).
\lstdefinestyle{chdeux}{language=Lean,
  literate={⁻}{{$^{-}$}}1 {⁻¹}{{$^{-1}$}}1 {ℝ}{{$\mathbb{R}$}}1 {ℤ}{{$\mathbb{Z}$}}1 {ℕ}{{$\mathbb{N}$}}1 {≤}{{$\le$}}1 {≥}{{$\ge$}}1 {≠}{{$\ne$}}1
  {→}{{$\to$}}1 {↔}{{$\leftrightarrow$}}1 {∀}{{$\forall$}}1 {∃}{{$\exists$}}1 {∧}{{$\wedge$}}1 {∨}{{$\vee$}}1
  {¬}{{$\neg$}}1 {⊤}{{$\top$}}1 {∈}{{$\in$}}1 {∉}{{$\notin$}}1 {∩}{{$\cap$}}1 {√}{{$\surd$}}1
  {π}{{$\pi$}}1 {ρ}{{$\rho$}}1 {κ}{{$\kappa$}}1 {α}{{$\alpha$}}1 {δ}{{$\delta$}}1 {ε}{{$\varepsilon$}}1 {σ}{{$\sigma$}}1
  {•}{{$\bullet$}}1 {∑}{{$\sum$}}1 {²}{{$^2$}}1 {₀}{{$_0$}}1 {₁}{{$_1$}}1 {₂}{{$_2$}}1 {‖}{{$\|$}}1 {⟨}{{$\langle$}}1 {⟩}{{$\rangle$}}1
  {𝓝}{{$\mathcal{N}$}}1 {∞}{{$\infty$}}1 {×}{{$\times$}}1 {·}{{$\cdot$}}1 {∘}{{$\circ$}}1 {∖}{{$\setminus$}}1
  {ψ}{{$\psi$}}1 {φ}{{$\varphi$}}1 {λ}{{$\lambda$}}1 {μ}{{$\mu$}}1 {θ}{{$\theta$}}1 {Λ}{{$\Lambda$}}1 {∂}{{$\partial$}}1
  {—}{{---}}1 {–}{{--}}1 {é}{{\'e}}1 {ö}{{\"o}}1 {’}{{'}}1
  {β}{{$\beta$}}1 {∫}{{$\int$}}1 {ᵐ}{{$^{\mathrm{m}}$}}1 {ι}{{$\iota$}}1 {ᵢ}{{$_i$}}1 {△}{{$\triangle$}}1
  {⊥}{{$\bot$}}1 {χ}{{$\chi$}}1}

\chapter{Le flux des formalisations : de l'optimalité universelle aux potentiels physiques}\label{ch2:chapitre}
\chaptermark{Le flux des formalisations}

\begin{center}
\begin{minipage}{0.9\linewidth}\small\itshape
Au chapitre~1, un théorème venu d'ailleurs --- le réseau triangulaire minimise l'énergie de tout potentiel complètement
monotone du carré de la distance --- a été posé, lu et vérifié par le noyau de Lean. Ce chapitre raconte ce qui s'est
passé ensuite : comment, en suivant une méthode fixée à l'avance, on a appliqué ce théorème à trois interactions de la
physique des matériaux, comment une petite bibliothèque a transformé ces applications en quelques lignes, et comment la
même méthode, tournée vers un problème voisin beaucoup plus difficile (la conjecture de Sandier--Serfaty), a surtout
servi à dire précisément ce qui \emph{n'est pas} démontré.
\end{minipage}
\end{center}

\begin{warnbox}[: à lire avant tout le chapitre]
\begin{itemize}[leftmargin=1.2em,itemsep=2pt]
\item Tout ce qui repose sur openai/math est \emph{conditionnel} aux théorèmes amont, qui sont produits par une IA et
\emph{non relus par des pairs}.
\item Chaque résultat de ce chapitre a été produit par des agents d'IA (Claude) sous la direction de l'auteur, puis
audité par un \emph{autre} agent d'IA. \textbf{Aucun humain n'a audité ces résultats}, ni l'adéquation des énoncés Lean,
ni les démonstrations « humaines » écrites ici.
\item Un \tierA{} porte sur l'énoncé Lean, jamais sur sa lecture physique.
\item La conjecture de Sandier--Serfaty \textbf{n'est pas démontrée} dans ce livre (section~\ref{ch2:sec:ss}).
\end{itemize}
\end{warnbox}

% ---------------------------------------------------------------------
\section{Vue d'ensemble : un flux, pas un catalogue}\label{ch2:sec:flux}

Le point de départ (chapitre~1) est le théorème amont \lean{universal_energy_minimum} d'openai/math : pour tout potentiel
$g$ \emph{admissible}, le réseau triangulaire $\Lat$ de covolume un a une énergie inférieure ou égale à celle de toute
configuration localement finie de densité un \tierA{}, conditionnellement à l'amont~\cite{oai}. Dans la
littérature humaine, l'optimalité universelle est démontrée en dimensions 8 et 24~\cite{ckmrv} \tierL, et
Petrache et Serfaty écrivaient en 2019 que le cas plan restait ouvert~\cite{ch2ps} \tierL : le théorème amont est donc
une revendication d'IA sur un problème ouvert, acceptée ici uniquement comme énoncé Lean vérifié par le noyau à partir
des définitions amont. Rappelons les deux énoncés, cités mot pour mot.

\begin{leanbox}[--- openai/math, \lean{OAI/Analysis/Triangular/Energy/Basic.lean:26-29}]
\begin{lstlisting}[style=chdeux]
def AdmissiblePotential (g : ℝ → ℝ) : Prop :=
  ContDiffOn ℝ (⊤ : ℕ∞) g (Set.Ioi 0) ∧
  (∀ t : ℝ, 0 < t → 0 ≤ g t) ∧
  ∀ (r : ℕ) (t : ℝ), 0 < t → 0 ≤ (-1 : ℝ) ^ r * iteratedDeriv r g t
\end{lstlisting}
\end{leanbox}
\begin{leanbox}[--- openai/math, \lean{OAI/Analysis/Triangular/Energy/Universal.lean:9-11} (preuve omise)]
\begin{lstlisting}[style=chdeux]
theorem universal_energy_minimum (g : ℝ → ℝ) (C : Set Plane)
    (hg : AdmissiblePotential g) (_hC : LocallyFinite C) (hd : DensityOne C) :
    latticeEnergy g ≤ energy g C ∧ latticeEnergy g = energy g A := by
\end{lstlisting}
\end{leanbox}
\noindent (Ces deux extraits sont cités d'après la copie épinglée d'openai/math ; la copie de travail
\lean{results/openai_math/hypotheses/night_2026-10-08/TRI_D3/src/} du dépôt lui est identique pour ces deux fichiers.) La variable de $g$ est le \emph{carré} $t=r^2$ de la
distance ; « admissible » veut dire : $C^\infty$ sur $]0,+\infty[$, positif, et $(-1)^r g^{(r)}(t)\ge0$ pour tout $r$,
c'est-à-dire \emph{complètement monotone} (CM) sur $]0,+\infty[$. Les énergies vivent dans $[0,+\infty]$.

Le flux de ce chapitre est le suivant (figure~\ref{ch2:fig:flux}).
\begin{enumerate}[itemsep=1pt]
\item \textbf{Amont} (chapitre~1) : l'optimalité universelle en dimension deux \tierA{} conditionnel, et nos corollaires
D1--D4 (Riesz, toute densité, sens CKMRV, disques ouverts ou fermés), publiés dans le manuel~v1
(doi\string:10.5281/zenodo.23295142)~\cite{v1}.
\item \textbf{Trois potentiels physiques} (section~\ref{ch2:sec:pot}) : excitons dipolaires, Yukawa plan, Rytova--Keldysh.
Pour chacun : une démonstration de complète monotonie en $r^2$, l'énoncé Lean, la conséquence.
\item \textbf{La généralisation} (section~\ref{ch2:sec:lib}) : une bibliothèque CM/Bernstein dont le théorème de
subordination fait de ces trois résultats des corollaires de deux ou trois lignes.
\item \textbf{Un problème voisin} (section~\ref{ch2:sec:ss}) : la conjecture de Sandier--Serfaty, où la même méthode
produit une reproduction, un audit, un énoncé Lean, une formalisation \emph{partielle} et un résultat négatif en
dimension trois --- mais aucune preuve.
\end{enumerate}

\begin{figure}[ht]\centering
\begin{tikzpicture}[x=1cm,y=1cm,font=\small,
  bx/.style={draw,rounded corners=2pt,align=center,inner sep=4pt,fill=#1!6},
  ar/.style={->,>=stealth,thick,gray!70!black}]
\node[bx=tierL] (up) at (0,0) {openai/math\\ \lean{universal_energy_minimum}\\ (IA, non relu)};
\node[bx=tierA] (d) at (4.6,0) {D1--D4\\ (chap.~1)};
\node[bx=tierA] (h1) at (0,-2.2) {H1 dipolaire\\ $2(t^{-1/2}-(t+d^2)^{-1/2})$};
\node[bx=tierA] (h2) at (4.6,-2.2) {H2 Yukawa\\ $e^{-\kappa\sqrt t}/\sqrt t$};
\node[bx=tierA] (h3) at (9.2,-2.2) {H3 Keldysh\\ mélange de Laplace};
\node[bx=tierA] (p3) at (4.6,-4.4) {P3 : bibliothèque CM/Bernstein\\ subordination $f\circ\psi$};
\node[bx=tierC] (ss) at (9.2,0) {Sandier--Serfaty\\ P1, A1, B3\\ (non démontrée)};
\draw[ar] (up) -- (d);
\draw[ar] (up) -- (h1);
\draw[ar] (up) -- (h2);
\draw[ar] (up.south east) -- (h3.north west);
\draw[ar] (h1) -- (p3);
\draw[ar] (h2) -- (p3);
\draw[ar] (h3) -- (p3);
\draw[ar,dashed] (d) -- node[above,font=\scriptsize]{même méthode} (ss);
\end{tikzpicture}
\caption{Le flux des formalisations de ce chapitre. Vert : énoncés Lean \tierA{} (conditionnels à l'amont) ;
bleu : amont cité ; orange : travail dont le statut mathématique reste \tierC{} ou inférieur.}\label{ch2:fig:flux}
\end{figure}

% ---------------------------------------------------------------------
\section{La méthode : préenregistrer, falsifier, laisser juger le noyau}\label{ch2:sec:methode}

Avant le moindre calcul, il faut dire comment on saura qu'on s'est trompé. La méthode suivie dans tous les travaux de ce
chapitre tient en quatre règles. Elle est outillée par Elenchus~\cite{elenchus}, un registre de revendications
(« ledger ») muni d'une porte automatique.

\subsection{Règle 1 : préenregistrer}\label{ch2:ssec:prereg}
On écrit, \emph{avant} de voir les résultats, ce qu'on cherche, ce qu'on rapportera et ce qui ferait échouer le
projet. Deux exemples réels.
\begin{itemize}
\item Pour les excitons dipolaires, le préenregistrement \texttt{DIP-TRI-1} a été validé (commit
\texttt{4feb027}) \emph{avant} le code du solveur (commit \texttt{17ea6be}) dans rusty-SUNDIALS~\cite{ch2dip}. Il fixait
la quantité à rapporter (une somme d'images tronquée plus la queue uniforme). L'article rapporte à la place des énergies
d'Ewald exactes, et le dit explicitement comme \emph{déviation} : la somme tronquée laissait un résidu non monotone.
\item Pour la formalisation du certificat de Sandier--Serfaty (A1), les probabilités a priori étaient écrites dans
\texttt{preregistration.json} : $0{,}6$ que quelques étapes pilotes aboutissent, $0{,}2$ que toutes aboutissent, $0{,}35$
pour un lemme complet de la réduction, $0{,}02$ pour l'ensemble ; et un \emph{critère d'arrêt} : aucune annonce de preuve
formelle tant que toutes les étapes et toute la réduction ne sont pas vérifiées par le noyau~\cite{ch2ssl}. Le critère
s'applique (section~\ref{ch2:ssec:a1}).
\end{itemize}

\subsection{Règle 2 : une porte de falsification, avec témoins positifs et négatifs}\label{ch2:ssec:falsif}
Toute expérience, tout théorème, tout programme est accompagné d'un \emph{témoin positif} (un cas où la réponse doit
être « oui ») et d'un \emph{témoin négatif} (un cas où elle doit être « non »), exécutés avant de rapporter le moindre
nombre. Une porte qui accepte tout ne vérifie rien. Dans ce chapitre, chaque énoncé d'admissibilité a son contre-exemple
vérifié par le noyau : $t^{-1/2}-2(t+d^2)^{-1/2}$ pour les dipôles, $e^{-\kappa\sqrt t}\cos t$ et $g_1-2g_2$ pour
Yukawa, $e^{-t}-2e^{-2t}$ pour les mélanges de Laplace, $\cos$ et une composition « poubelle » pour la bibliothèque.
Pour le certificat de Sandier--Serfaty, trois versions volontairement affaiblies du programme ont été exécutées ; la
section~\ref{ch2:ssec:p1} dit honnêtement que l'une d'elles \emph{passe}.

\subsection{Règle 3 : le noyau est le juge --- mais pas son code de retour}\label{ch2:ssec:noyau}
Un fichier Lean qui contient \lean{sorry}, ou un axiome planté, \emph{compile avec le code de retour 0}. Le code de
retour n'est donc pas la porte. Voici le fichier de témoins plantés de la bibliothèque CM, cité mot pour mot.
\begin{leanbox}[--- \lean{formal_cert/Programmes/CompletelyMonotone/Controls.lean:11-18}]
\begin{lstlisting}[style=chdeux]
theorem ProgCM.ctl_sorry : ProgCM.IsCompletelyMonotoneOn Real.cos := sorry

axiom ProgCM.ctl_cheat : False

theorem ProgCM.ctl_axiom : ProgCM.IsCompletelyMonotoneOn Real.cos := ProgCM.ctl_cheat.elim

#print axioms ProgCM.ctl_sorry
#print axioms ProgCM.ctl_axiom
\end{lstlisting}
\end{leanbox}
\noindent Les deux « preuves » de l'énoncé faux « $\cos$ est CM » compilent avec le code~0 ; seule la commande
\texttt{\#print axioms} les démasque (\lean{sorryAx}, resp.\ \lean{ProgCM.ctl_cheat})~\cite{p3}. D'où la définition du
\tierA{} (fiche \texttt{FACTS.md}) : la liste des axiomes est contenue dans $\{$\lean{propext}, \lean{Classical.choice},
\lean{Quot.sound}$\}$, plus un rejeu par le noyau (\texttt{leanchecker}), plus une recherche textuelle de
\lean{skipKernelTC}, \lean{addDecl}, \lean{native_decide}, \lean{implemented_by}.

\begin{remark}[Pourquoi \texttt{\#print axioms} ne suffit pas]\label{ch2:rem:skip}
L'audit A2 du corpus Navier--Stokes d'openai a montré qu'une déclaration forgée avec l'option
\lean{debug.skipKernelTC} peut afficher une liste d'axiomes irréprochable. Le rejeu complet par le noyau ferme cette
porte dérobée. Par honnêteté : les articles sources de ces sections ne rapportaient pas de rejeu \texttt{leanchecker}. Ce rejeu a été
fait pour ce livre le 2026-10-26 (enregistrement \texttt{results/kernel\_replay\_2026-10-26/}), sur les modules
cités ici : code de retour~0 pour chacun des modules dipolaires (\lean{DipolarExciton}, \lean{DipolarExcitonMultipolar},
\lean{DipolarExcitonStep1}), de Yukawa (\lean{Yukawa}, \lean{YukawaMix}), de Keldysh (\lean{KeldyshCM}) et de la
bibliothèque CM/Bernstein (\lean{Basic}, \lean{Bernstein}, \lean{Laplace}, \lean{Corollaries}, \lean{Applications},
\lean{Examples}, \lean{Limits}), et la recherche textuelle ne trouve aucun jeton interdit hors des fichiers de contrôles
plantés. Le rejeu porte sur les déclarations de ces modules ; Mathlib et le cœur de Lean sont chargés, non rejoués.
Les modules de la section~\ref{ch2:sec:ss} (Sandier--Serfaty) figurent dans le plan de rejeu mais n'avaient pas été
rejoués à l'assemblage de ce livre : leurs énoncés portent l'étiquette \tierAs{}.
\end{remark}

\subsection{Règle 4 : le producteur n'est pas le vérificateur}\label{ch2:ssec:prodverif}
L'agent qui écrit un théorème n'est jamais celui qui le déclare vérifié. Un second agent d'IA recompile les fichiers
\emph{depuis les sources}, dans son propre répertoire, plante ses propres témoins, relit les énoncés de manière
adversariale, et rend un verdict (par exemple \texttt{ACCEPT\_WITH\_CORRECTIONS}). Le registre Elenchus attribue à chaque
revendication un niveau --- A (noyau), L (littérature lue et épinglée), C (argument humain non vérifié par machine),
X (exploration en virgule flottante, qui ne soutient jamais une revendication) --- et la porte vérifie qu'aucune ligne
n'est classée au-dessus de ce sur quoi elle repose, transitivement. Elle lève un drapeau
\texttt{LEDGER\_UNAUDITED\_TIER\_A} pour chaque ligne de niveau A dont aucun humain n'a certifié que l'énoncé veut dire ce
que dit sa glose. Ces drapeaux restent ouverts dans tous les travaux de ce chapitre : un audit par IA ne les ferme pas.

\begin{warnbox}[: ce que certifie un niveau A]
Le noyau certifie qu'\emph{un énoncé Lean} est démontré à partir des définitions Lean. Il ne certifie ni que ces
définitions traduisent fidèlement la physique, ni que l'amont (openai/math) est correct au sens humain : il certifie
que, \emph{si} l'on accepte les définitions et les théorèmes amont, la conclusion suit.
\end{warnbox}

% ---------------------------------------------------------------------
\section{Trois potentiels physiques, trois démonstrations}\label{ch2:sec:pot}

Les trois sous-sections suivent le même plan : (i) le modèle physique \tierC{} ; (ii) une démonstration pas à pas que
le potentiel, vu comme fonction de $t=r^2$, est complètement monotone \tierC{}, avec en regard l'énoncé Lean vérifié
\tierA{} ; (iii) la conséquence : optimalité du réseau triangulaire (conditionnelle), finitude de l'énergie quand elle
est démontrée ; (iv) un témoin négatif.

\begin{definition}[complète monotonie]\label{ch2:def:cm}
Une fonction $f:\,]0,+\infty[\,\to\R$ est \emph{complètement monotone} (CM) si elle est de classe $C^\infty$ et si, pour
tout $n\in\N$ et tout $t>0$, $(-1)^nf^{(n)}(t)\ge0$. Autrement dit : $f\ge0$, $f'\le0$, $f''\ge0$, $f'''\le0$, etc.
\end{definition}

\begin{example}
$t\mapsto e^{-at}$ ($a\ge0$) est CM car $(-1)^n\frac{d^n}{dt^n}e^{-at}=a^ne^{-at}\ge0$. La fonction $\cos$ ne l'est pas :
$\cos\pi=-1<0$. La fonction $\sqrt t$ non plus : elle est croissante.
\end{example}

\begin{remark}[pourquoi la variable $t=r^2$ ?]
Le théorème amont parle de $g(\|x-y\|^2)$. Être CM en $r^2$ est plus exigeant qu'être CM en $r$ : $e^{-r}$ est CM en $r$,
et l'on verra que $e^{-\sqrt t}$ est aussi CM en $t$, mais cela demande un argument (section~\ref{ch2:ssec:yuk}).
\end{remark}

% .....................................................................
\subsection{Excitons dipolaires dans une bicouche}\label{ch2:ssec:dip}

\begin{physbox}[: le modèle \tierC]
Dans une hétérobicouche de dichalcogénures de métaux de transition (par exemple MoSe$_2$/WSe$_2$), un exciton
« indirect » est un électron dans une couche, à la hauteur $d$, lié à un trou dans l'autre couche, à la hauteur $0$,
à la même position dans le plan. Deux excitons à distance $r$ dans le plan forment quatre paires de Coulomb : électron--électron
et trou--trou à distance $r$ (répulsion $+1/r$ chacune), et deux paires croisées à distance $\sqrt{r^2+d^2}$ (attraction
$-1/\sqrt{r^2+d^2}$ chacune). On néglige l'énergie de liaison interne, constante. En unités $e^2/(4\pi\varepsilon)$ :
\[V(r)=\frac2r-\frac2{\sqrt{r^2+d^2}},\qquad g_d(t)=2\bigl(t^{-1/2}-(t+d^2)^{-1/2}\bigr),\quad V(r)=g_d(r^2).\]
À grande distance $V(r)=d^2/r^3+O(r^{-5})$ : la loi dipôle--dipôle. Ce comptage de paires est \tierC{} ; l'identité
algébrique $g_d(r^2)=2(1/r-1/\sqrt{r^2+d^2})$ est vérifiée par le noyau (\lean{g_inplane}, \tierA{} : le module \texttt{DipolarExcitonStep1.lean} n'était vérifié que par son producteur dans l'article source~\cite{ch2dip} ; il a depuis été rejoué par \lean{leanchecker}, code~0).
\end{physbox}

\begin{leanbox}[--- \lean{formal_cert/Autoresearch/DipolarExciton.lean:22}]
\begin{lstlisting}[style=chdeux]
noncomputable def g (d : ℝ) : ℝ → ℝ := fun t => 2 * (t ^ (-(1 / 2 : ℝ)) - (t + d ^ 2) ^ (-(1 / 2 : ℝ)))
\end{lstlisting}
\end{leanbox}

\begin{theorem}[H1, admissibilité du potentiel dipolaire]\label{ch2:thm:dip}
Pour tout réel $d$, $g_d$ est complètement monotone sur $]0,+\infty[$. Plus généralement, pour $s>0$ et $c\ge0$,
$\mathrm{shiftDiff}_{s,c}(t)=t^{-s/2}-(t+c)^{-s/2}$ est CM. \tierA{} (énoncé Lean) ; démonstration ci-dessous \tierC.
\end{theorem}

La démonstration repose sur les \emph{symboles de Pochhammer}. Pour $a\in\R$ et $r\in\N$, on note
\[\mathrm{desc}_r(a)=a(a-1)\cdots(a-r+1),\qquad \mathrm{asc}_r(a)=a(a+1)\cdots(a+r-1),\]
avec $\mathrm{desc}_0=\mathrm{asc}_0=1$ (factorielles descendante et montante ; Mathlib : \lean{descPochhammer},
\lean{ascPochhammer}).

\begin{lemma}[dérivées d'une puissance décalée]\label{ch2:lem:pochder}
Pour $a\in\R$, $c\ge0$, $t>0$ et $r\in\N$ :
$\displaystyle\frac{d^r}{dt^r}(t+c)^a=\mathrm{desc}_r(a)\,(t+c)^{a-r}$.
\end{lemma}
\begin{humanproof}
Récurrence sur $r$. Pour $r=0$ c'est l'identité. Si la formule vaut au rang $r$, on dérive $(t+c)^{a-r}$, ce qui est
licite car $t+c>0$ : on obtient $(a-r)(t+c)^{a-r-1}$, et $\mathrm{desc}_r(a)(a-r)=\mathrm{desc}_{r+1}(a)$.
\end{humanproof}
\begin{leanbox}[--- \lean{DipolarExciton.lean:32-33} et \lean{:36-37} (preuves omises)]
\begin{lstlisting}[style=chdeux]
lemma iteratedDeriv_rpow (a : ℝ) (r : ℕ) (t : ℝ) :
    iteratedDeriv r (fun x : ℝ => x ^ a) t = (descPochhammer ℝ r).eval a * t ^ (a - r) := by
\end{lstlisting}
\begin{lstlisting}[style=chdeux]
lemma iteratedDeriv_shift (a c : ℝ) (r : ℕ) (t : ℝ) :
    iteratedDeriv r (fun x : ℝ => (x + c) ^ a) t = (descPochhammer ℝ r).eval a * (t + c) ^ (a - r) := by
\end{lstlisting}
\end{leanbox}

Par linéarité de la dérivation, pour $t>0$ :
\begin{equation}\label{ch2:eq:gder}
g_d^{(r)}(t)=2\,\mathrm{desc}_r\!\left(-\tfrac12\right)\Bigl(t^{-1/2-r}-(t+d^2)^{-1/2-r}\Bigr).
\end{equation}
\begin{leanbox}[--- \lean{DipolarExciton.lean:41-44} (preuve omise)]
\begin{lstlisting}[style=chdeux]
lemma iteratedDeriv_g (d : ℝ) (r : ℕ) {t : ℝ} (ht : 0 < t) :
    iteratedDeriv r (g d) t =
      2 * ((descPochhammer ℝ r).eval (-(1 / 2 : ℝ)) *
        (t ^ (-(1 / 2 : ℝ) - r) - (t + d ^ 2) ^ (-(1 / 2 : ℝ) - r))) := by
\end{lstlisting}
\end{leanbox}

\begin{lemma}[le signe]\label{ch2:lem:sign}
$(-1)^r\,\mathrm{desc}_r(-\tfrac12)=\mathrm{asc}_r(\tfrac12)=\tfrac12\cdot\tfrac32\cdots\left(r-\tfrac12\right)>0$.
\end{lemma}
\begin{humanproof}
$(-1)^r\mathrm{desc}_r(-\tfrac12)=\prod_{k=0}^{r-1}(-1)\left(-\tfrac12-k\right)=\prod_{k=0}^{r-1}\left(\tfrac12+k\right)$,
produit de nombres strictement positifs.
\end{humanproof}
\begin{leanbox}[--- \lean{DipolarExciton.lean:56-57} (preuve omise)]
\begin{lstlisting}[style=chdeux]
lemma sign_desc (r : ℕ) :
    (-1 : ℝ) ^ r * (descPochhammer ℝ r).eval (-(1 / 2 : ℝ)) = (ascPochhammer ℝ r).eval (1 / 2 : ℝ) := by
\end{lstlisting}
\end{leanbox}

\begin{lemma}[le crochet]\label{ch2:lem:crochet}
Pour $c\ge0$, $e\le0$ et $t>0$ : $t^e-(t+c)^e\ge0$.
\end{lemma}
\begin{humanproof}
Pour $e\le0$, $x\mapsto x^e$ est décroissante sur $]0,+\infty[$ ; or $t\le t+c$.
\end{humanproof}
\begin{leanbox}[--- \lean{DipolarExciton.lean:62-63} (preuve omise)]
\begin{lstlisting}[style=chdeux]
lemma diff_nonneg (c e : ℝ) (hc : 0 ≤ c) (he : e ≤ 0) {t : ℝ} (ht : 0 < t) :
    0 ≤ t ^ e - (t + c) ^ e := by
\end{lstlisting}
\end{leanbox}

\begin{humanproof}[Démonstration du théorème~\ref{ch2:thm:dip}]
\emph{Régularité.} $t\mapsto t^{-1/2}$ et $t\mapsto(t+d^2)^{-1/2}$ sont $C^\infty$ là où leur base est strictement
positive, donc sur $]0,+\infty[$.
\emph{Signes.} D'après~\eqref{ch2:eq:gder} et le lemme~\ref{ch2:lem:sign},
\[(-1)^rg_d^{(r)}(t)=2\,\mathrm{asc}_r\!\left(\tfrac12\right)\Bigl(t^{-1/2-r}-(t+d^2)^{-1/2-r}\Bigr),\]
produit d'un nombre positif par un crochet positif (lemme~\ref{ch2:lem:crochet} avec $e=-\tfrac12-r\le0$ et
$c=d^2\ge0$). Le cas $r=0$ donne $g_d\ge0$. En remplaçant $\tfrac12$ par $\tfrac s2$ ($s>0$) et $d^2$ par $c\ge0$, rien
ne change : c'est l'énoncé pour $\mathrm{shiftDiff}_{s,c}$. Enfin $g_d=2\,\mathrm{shiftDiff}_{1,d^2}$ et un multiple
positif d'une fonction CM est CM, ce qui donne le résultat pour \emph{tout} réel $d$.
\end{humanproof}

\begin{leanbox}[--- \lean{DipolarExciton.lean:68} et \lean{DipolarExcitonMultipolar.lean:37, 59, 102} (preuves omises)]
\begin{lstlisting}[style=chdeux]
theorem dipolar_admissible {d : ℝ} (_hd : 0 < d) : AdmissiblePotential (g d) := by
\end{lstlisting}
\begin{lstlisting}[style=chdeux]
noncomputable def shiftDiff (s c : ℝ) : ℝ → ℝ := fun t => t ^ (-(s / 2)) - (t + c) ^ (-(s / 2))
\end{lstlisting}
\begin{lstlisting}[style=chdeux]
theorem shiftDiff_admissible {s c : ℝ} (hs : 0 < s) (hc : 0 ≤ c) : AdmissiblePotential (shiftDiff s c) := by
\end{lstlisting}
\begin{lstlisting}[style=chdeux]
theorem dipolar_admissible_all (d : ℝ) : AdmissiblePotential (AutoresearchDipolar.g d) := by
\end{lstlisting}
\end{leanbox}

\begin{remark}[lecture de Bernstein \tierC]
De façon équivalente, $t^{-1/2}-(t+c)^{-1/2}=\pi^{-1/2}\int_0^\infty e^{-ut}\,(1-e^{-uc})\,u^{-1/2}\,du$ est un mélange
positif d'exponentielles, ce qui rend la complète monotonie évidente~\cite{ch2dip}. Cette identité n'est pas vérifiée
en Lean ; la preuve par Pochhammer ci-dessus l'évite.
\end{remark}

\paragraph{Conséquence 1 : optimalité à densité un \tierA, conditionnelle.} On applique le théorème amont à $g_d$.
La preuve Lean est d'une ligne : la première composante de \lean{universal_energy_minimum}.
\begin{leanbox}[--- \lean{DipolarExciton.lean:94-96} (énoncé et preuve complète)]
\begin{lstlisting}[style=chdeux]
theorem triangular_dipolar_optimal {d : ℝ} (hd : 0 < d) (C : Set Plane) (hC : LocallyFinite C)
    (hdens : DensityOne C) : latticeEnergy (g d) ≤ energy (g d) C :=
  (universal_energy_minimum (g d) C (dipolar_admissible hd) hC hdens).1
\end{lstlisting}
\end{leanbox}

\paragraph{Conséquence 2 : l'inégalité n'est pas $\top\le\top$ \tierA.} Une inégalité dans $[0,+\infty]$ dont le membre
de gauche vaut $+\infty$ ne dit rien. Il faut donc montrer que $\mathrm{latticeEnergy}(g_d)<+\infty$.
\begin{lemma}\label{ch2:lem:diffle}
Pour $c\ge0$ et $t>0$ : $t^{-1/2}-(t+c)^{-1/2}\le\frac c2\,t^{-3/2}$. En particulier $g_d(t)\le d^2t^{-3/2}$.
\end{lemma}
\begin{humanproof}
Posons $a=t^{-1/2}$ et $b=(t+c)^{-1/2}$, de sorte que $0<b\le a$, $a^2t=1$ et $b^2(t+c)=1$. Alors
$a^2-b^2=a^2b^2\bigl((t+c)-t\bigr)=c\,a^2b^2$, soit $(a-b)(a+b)=c\,a^2b^2$. Comme $b\le a$, $2b^2\le a^2+ab=a(a+b)$, donc
$2(a-b)(a+b)\le c\,a^3(a+b)$ ; en divisant par $a+b>0$ : $a-b\le\frac c2a^3=\frac c2t^{-3/2}$.
\end{humanproof}
\noindent La somme $\sum_{a\in\Lat\setminus\{0\}}\|a\|^{-3}$ converge (exposant de Riesz $3>2$, chapitre~1), donc
$\mathrm{latticeEnergy}(g_d)\le d^2\,\mathrm{latticeEnergy}(t^{-3/2})<+\infty$.
\begin{leanbox}[--- \lean{DipolarExciton.lean:119-120, 159, 167} (preuves omises)]
\begin{lstlisting}[style=chdeux]
lemma diff_le (c : ℝ) (hc : 0 ≤ c) {t : ℝ} (ht : 0 < t) :
    t ^ (-(1 / 2 : ℝ)) - (t + c) ^ (-(1 / 2 : ℝ)) ≤ c / 2 * t ^ (-(3 / 2 : ℝ)) := by
\end{lstlisting}
\begin{lstlisting}[style=chdeux]
lemma g_le (d : ℝ) {t : ℝ} (ht : 0 < t) : g d t ≤ d ^ 2 * riesz 3 t := by
\end{lstlisting}
\begin{lstlisting}[style=chdeux]
theorem latticeEnergy_dipolar_lt_top (d : ℝ) : latticeEnergy (g d) < ⊤ := by
\end{lstlisting}
\end{leanbox}

\paragraph{Conséquence 3 : toute densité $\rho>0$ \tierA, conditionnelle.} Si $C$ a la densité $\rho$, alors
$\sqrt\rho\,C$ a la densité un et $\|\sqrt\rho x-\sqrt\rho y\|^2=\rho\|x-y\|^2$ : l'énergie de $C$ pour $g$ est celle de
$\sqrt\rho\,C$ pour $t\mapsto g(t/\rho)$, qui reste CM (chaque dérivée gagne le facteur $\rho^{-r}>0$). Le corollaire D2
du chapitre~1 donne alors l'optimalité de $\rho^{-1/2}\Lat$, et $g_d(t/\rho)\le d^2\rho^{3/2}t^{-3/2}$ donne la finitude.
\begin{leanbox}[--- \lean{DipolarExcitonMultipolar.lean:125-129} (énoncé et preuve complète) et \lean{:175-176}]
\begin{lstlisting}[style=chdeux]
theorem triangular_dipolar_optimal_any_density {d : ℝ} (hd : 0 < d) {ρ : ℝ} (hρ : 0 < ρ) :
    LocallyFinite ((Real.sqrt ρ)⁻¹ • A) ∧ TriangularDensity.DensityRho ρ ((Real.sqrt ρ)⁻¹ • A) ∧
      ∀ C : Set Plane, LocallyFinite C → TriangularDensity.DensityRho ρ C →
        energy (AutoresearchDipolar.g d) ((Real.sqrt ρ)⁻¹ • A) ≤ energy (AutoresearchDipolar.g d) C :=
  TriangularDensity.triangular_optimal_any_density _ (AutoresearchDipolar.dipolar_admissible hd) hρ
\end{lstlisting}
\begin{lstlisting}[style=chdeux]
theorem energy_dipolar_any_density_lt_top {d : ℝ} (hd : 0 < d) {ρ : ℝ} (hρ : 0 < ρ) :
    energy (AutoresearchDipolar.g d) ((Real.sqrt ρ)⁻¹ • A) < ⊤ := by
\end{lstlisting}
\end{leanbox}

\paragraph{Témoins négatifs \tierA.} (i) Si l'on double le poids de l'attraction croisée,
$h_d(t)=t^{-1/2}-2(t+d^2)^{-1/2}$, alors en $t=d^2$ :
$h_d(d^2)=\frac1d-\frac{2}{\sqrt2\,d}<\frac1d-\frac{2}{2d}=0$, donc $h_d$ n'est pas admissible (\lean{h_not_admissible}).
(ii) Pour $c<0$ et $t=-2c>0$, on a $t+c=-c\in\,]0,t[$, d'où $(t+c)^{-s/2}>t^{-s/2}$ et
$\mathrm{shiftDiff}_{s,c}(t)<0$ (\lean{shiftDiff_not_admissible_of_neg}). Le signe du décalage est porteur.

\begin{warnbox}[: portée]
La finitude est démontrée pour $g_d$ seulement ; pour la famille $\mathrm{shiftDiff}_{s,c}$ aucune finitude n'est
démontrée, et pour $c=0$ on a $\mathrm{shiftDiff}_{s,0}\equiv0$ : la conclusion est alors le trivial
$0\le\mathrm{energy}$. Le modèle est classique : dipôles rigides, pas d'énergie cinétique, pas d'écrantage diélectrique,
pas de potentiel de moiré. Dans la fenêtre de densité expérimentale, les excitons réels forment un gaz quantique, pas un
cristal classique \tierC~\cite{ch2dip}. Article : doi\string:10.5281/zenodo.23646103.
\end{warnbox}

% .....................................................................
\subsection{Le potentiel de Yukawa plan}\label{ch2:ssec:yuk}

\begin{physbox}[: le modèle \tierC]
Des particules chargées confinées dans un plan mais plongées dans un milieu écrantant à trois dimensions (plasmas
poussiéreux, colloïdes chargés) interagissent par le potentiel de Coulomb écranté $V(r)=e^{-\kappa r}/r$, $\kappa>0$.
En variable $t=r^2$ : $g_\kappa(t)=e^{-\kappa\sqrt t}/\sqrt t$. \textbf{Ce n'est pas} la fonction de Green
$K_0(\kappa r)$ de l'équation de Debye--Hückel bidimensionnelle, qui n'est pas traitée ici~\cite{ch2yuk}.
\end{physbox}
\begin{leanbox}[--- \lean{formal_cert/Autoresearch/Yukawa.lean:28, 31}]
\begin{lstlisting}[style=chdeux]
noncomputable def phi (κ : ℝ) : ℝ → ℝ := fun t => Real.exp (-κ * Real.sqrt t)
\end{lstlisting}
\begin{lstlisting}[style=chdeux]
noncomputable def g (κ : ℝ) : ℝ → ℝ := fun t => Real.exp (-κ * Real.sqrt t) / Real.sqrt t
\end{lstlisting}
\end{leanbox}

Le plan est : $g_\kappa=t^{-1/2}\cdot\varphi_\kappa$ avec $\varphi_\kappa(t)=e^{-\kappa\sqrt t}$ ; on montre que le
produit de deux fonctions CM est CM, puis que $\varphi_\kappa$ est CM par une récurrence forte. On dit que $f$ a le
\emph{motif de signes jusqu'à l'ordre $n$} en $t$ si $(-1)^if^{(i)}(t)\ge0$ pour $i=0,\dots,n$.

\begin{lemma}[signes de Leibniz]\label{ch2:lem:leibniz}
Si $f$ et $h$ sont $n$ fois dérivables en $t$ et ont toutes deux le motif de signes jusqu'à l'ordre $n$ en $t$, alors
$(-1)^n(fh)^{(n)}(t)\ge0$. En conséquence, le produit de deux fonctions CM est CM.
\end{lemma}
\begin{humanproof}
Formule de Leibniz : $(fh)^{(n)}=\sum_{i=0}^n\binom ni f^{(i)}h^{(n-i)}$. On écrit $(-1)^n=(-1)^i(-1)^{n-i}$ :
\[(-1)^n(fh)^{(n)}=\sum_{i=0}^n\binom ni\,\bigl[(-1)^if^{(i)}\bigr]\,\bigl[(-1)^{n-i}h^{(n-i)}\bigr],\]
somme de produits de nombres positifs. La régularité et la positivité d'un produit sont immédiates.
\end{humanproof}
\begin{leanbox}[--- \lean{Yukawa.lean:38-41} (preuve omise ; elle tient en une douzaine de lignes)]
\begin{lstlisting}[style=chdeux]
lemma mul_sign {f h : ℝ → ℝ} {t : ℝ} (r : ℕ) (hf : ContDiffAt ℝ r f t) (hh : ContDiffAt ℝ r h t)
    (Hf : ∀ i ≤ r, 0 ≤ (-1 : ℝ) ^ i * iteratedDeriv i f t)
    (Hh : ∀ i ≤ r, 0 ≤ (-1 : ℝ) ^ i * iteratedDeriv i h t) :
    0 ≤ (-1 : ℝ) ^ r * iteratedDeriv r (f * h) t := by
\end{lstlisting}
\end{leanbox}

\begin{lemma}[$t^{-s/2}$ est CM]\label{ch2:lem:riesz}
Pour $s>0$, $p(t)=t^{-s/2}$ vérifie $(-1)^np^{(n)}(t)=\mathrm{asc}_n(\tfrac s2)\,t^{-s/2-n}>0$.
\end{lemma}
\begin{humanproof}
C'est le lemme~\ref{ch2:lem:pochder} avec $c=0$ et $a=-\tfrac s2$, suivi du calcul de signe du lemme~\ref{ch2:lem:sign}
avec $\tfrac s2$ au lieu de $\tfrac12$. (Lean : \lean{TriangularRiesz.riesz_admissible}, chapitre~1.)
\end{humanproof}

\begin{proposition}[$e^{-\kappa\sqrt t}$ est CM]\label{ch2:prop:phi}
Pour $\kappa>0$, $\varphi_\kappa(t)=e^{-\kappa\sqrt t}$ est complètement monotone sur $]0,+\infty[$.
\end{proposition}
\begin{humanproof}
Par dérivation composée, sur $]0,+\infty[$ :
\[\varphi_\kappa'(t)=-\frac{\kappa}{2\sqrt t}\,e^{-\kappa\sqrt t}=-\frac\kappa2\,t^{-1/2}\,\varphi_\kappa(t).\tag{$\ast$}\]
On démontre par \emph{récurrence forte} sur $n$ la propriété $\mathcal P(n)$ : « pour tout $m\le n$ et tout $t>0$,
$(-1)^m\varphi_\kappa^{(m)}(t)\ge0$ ».

$\mathcal P(0)$ : $\varphi_\kappa>0$.

$\mathcal P(n)\Rightarrow\mathcal P(n+1)$ : seul l'ordre $m=n+1$ est nouveau. Les deux membres de $(\ast)$ sont égaux sur
l'ouvert $]0,+\infty[$, donc leurs dérivées de tous ordres aussi : $\varphi_\kappa^{(n+1)}=(\varphi_\kappa')^{(n)}
=-\frac\kappa2\,(t^{-1/2}\varphi_\kappa)^{(n)}$. Par le lemme~\ref{ch2:lem:riesz}, $t^{-1/2}$ a le motif de signes à tout
ordre ; par $\mathcal P(n)$, $\varphi_\kappa$ l'a jusqu'à l'ordre $n$. Le lemme~\ref{ch2:lem:leibniz} donne
$(-1)^n(t^{-1/2}\varphi_\kappa)^{(n)}\ge0$, d'où
\[(-1)^{n+1}\varphi_\kappa^{(n+1)}=\frac\kappa2\,(-1)^n(t^{-1/2}\varphi_\kappa)^{(n)}\ge0,\]
car $\kappa>0$. Le passage de $n$ à $n+1$ n'utilise les signes de $\varphi_\kappa$ que jusqu'à l'ordre $n$ : c'est
pourquoi la récurrence forte suffit et qu'on n'a besoin d'aucune formule de Faà di Bruno. La régularité vient de ce que
$\exp$ et $\sqrt{\,\cdot\,}$ sont $C^\infty$ sur $]0,+\infty[$.
\end{humanproof}
\begin{leanbox}[--- \lean{Yukawa.lean:82-83, 94-95, 121} (preuves omises)]
\begin{lstlisting}[style=chdeux]
lemma phi_hasDerivAt (κ : ℝ) {t : ℝ} (ht : 0 < t) :
    HasDerivAt (phi κ) (-(κ / 2) * (riesz 1 t * phi κ t)) t := by
\end{lstlisting}
\begin{lstlisting}[style=chdeux]
lemma phi_sign {κ : ℝ} (hκ : 0 < κ) :
    ∀ n : ℕ, ∀ m ≤ n, ∀ t : ℝ, 0 < t → 0 ≤ (-1 : ℝ) ^ m * iteratedDeriv m (phi κ) t := by
\end{lstlisting}
\begin{lstlisting}[style=chdeux]
theorem phi_admissible {κ : ℝ} (hκ : 0 < κ) : AdmissiblePotential (phi κ) :=
\end{lstlisting}
\end{leanbox}
\noindent Remarquez la forme de \lean{phi_sign} : la quantification « $\forall\,m\le n$ » est exactement la récurrence
forte de la démonstration humaine.

\begin{theorem}[H2, admissibilité de Yukawa]\label{ch2:thm:yuk}
Pour $\kappa>0$, $g_\kappa$ est admissible. \tierA
\end{theorem}
\begin{humanproof}
Pour $t>0$, $g_\kappa(t)=t^{-1/2}\varphi_\kappa(t)$ : produit de deux fonctions CM (lemmes~\ref{ch2:lem:riesz},
\ref{ch2:lem:leibniz} et proposition~\ref{ch2:prop:phi}). En Lean, $g_\kappa(0)=0$ par la convention $x/0=0$ ; mais
l'admissibilité ne regarde que $]0,+\infty[$, et deux fonctions égales sur un ouvert y ont les mêmes dérivées
(\lean{admissible_congr}).
\end{humanproof}
\begin{leanbox}[--- \lean{Yukawa.lean:133-134} (énoncé et preuve complète)]
\begin{lstlisting}[style=chdeux]
theorem yukawa_admissible {κ : ℝ} (hκ : 0 < κ) : AdmissiblePotential (g κ) :=
  admissible_congr (admissible_mul (riesz_admissible (s := 1) one_pos) (phi_admissible hκ)) (g_eqOn κ)
\end{lstlisting}
\end{leanbox}

\paragraph{Conséquence : optimalité et énergie finie \tierA, conditionnelle.} L'optimalité à densité un
(\lean{triangular_yukawa_optimal}) s'obtient comme pour les dipôles. Pour la finitude : pour $x\ge0$,
$e^x\ge x^2/2$ ; avec $x=\kappa\sqrt t$,
\[g_\kappa(t)=\frac1{e^{\kappa\sqrt t}\sqrt t}\le\frac{2}{\kappa^2\,t\sqrt t}=\frac2{\kappa^2}\,t^{-3/2},\]
et la somme de Riesz d'exposant $3$ converge. L'hypothèse $\kappa>0$ sert ici : pour $\kappa=0$, $g_0=t^{-1/2}$ est le
Riesz d'exposant $1$, dont la somme sur le réseau diverge.
\begin{leanbox}[--- \lean{Yukawa.lean:160, 175} (preuves omises)]
\begin{lstlisting}[style=chdeux]
lemma g_le {κ : ℝ} (hκ : 0 < κ) {t : ℝ} (ht : 0 < t) : g κ t ≤ 2 / κ ^ 2 * riesz 3 t := by
\end{lstlisting}
\begin{lstlisting}[style=chdeux]
theorem latticeEnergy_yukawa_lt_top {κ : ℝ} (hκ : 0 < κ) : latticeEnergy (g κ) < ⊤ := by
\end{lstlisting}
\end{leanbox}

\paragraph{Mélanges finis \tierA.} Les fonctions admissibles forment un cône convexe (les dérivées sont linéaires),
donc tout mélange $\sum_iw_ig_{\kappa_i}$ avec $w_i\ge0$, $\kappa_i>0$ est admissible, d'énergie finie car majoré par
$\bigl(\sum_iw_i\,2/\kappa_i^2\bigr)t^{-3/2}$ ; le résultat vaut aussi à toute densité (\lean{YukawaMix.lean}).
\begin{leanbox}[--- \lean{formal_cert/Autoresearch/YukawaMix.lean:129-130, 138-139} (preuve omise)]
\begin{lstlisting}[style=chdeux]
noncomputable def yukawaMix {ι : Type*} (s : Finset ι) (w κ : ι → ℝ) : ℝ → ℝ :=
  fun t => ∑ i ∈ s, w i * g (κ i) t
\end{lstlisting}
\begin{lstlisting}[style=chdeux]
theorem yukawaMix_admissible {ι : Type*} (s : Finset ι) {w κ : ι → ℝ} (hw : ∀ i ∈ s, 0 ≤ w i)
    (hκ : ∀ i ∈ s, 0 < κ i) : AdmissiblePotential (yukawaMix s w κ) := by
\end{lstlisting}
\end{leanbox}

\paragraph{Témoins négatifs \tierA.} (a) $c_\kappa(t)=e^{-\kappa\sqrt t}\cos t$ prend la valeur
$-e^{-\kappa\sqrt\pi}<0$ en $t=\pi$ (\lean{c_not_admissible}). (b) Le mélange signé $g_1-2g_2$ : en $t=\frac14$,
$\sqrt t=\frac12$ et $g_1(\tfrac14)-2g_2(\tfrac14)=2e^{-1/2}-4e^{-1}=2e^{-1}(e^{1/2}-2)<0$ car $e<4$
(\lean{signedMix_not_admissible}). Ces témoins montrent que les hypothèses des lemmes ne sont pas vides ; ils ne
montrent \emph{pas} que le réseau triangulaire cesse d'être optimal pour ces fonctions : rien ici ne le décide.

\begin{warnbox}[: portée]
Modèle classique de particules ponctuelles ; aucune énergie cinétique ; pas le noyau $K_0$. L'article compagnon
(doi\string:10.5281/zenodo.23646105) contient en outre une étude numérique \tierX{} par CVODE qui corrige un erratum de
troncature publié ; elle ne soutient aucun énoncé de ce chapitre~\cite{ch2yuk}.
\end{warnbox}

% .....................................................................
\subsection{Le potentiel de Rytova--Keldysh}\label{ch2:ssec:kel}

\begin{physbox}[: le modèle \tierL{} / \tierC]
Dans une monocouche de semi-conducteur bidimensionnel (MoS$_2$, WSe$_2$, \dots) suspendue dans le vide, deux charges
interagissent par le potentiel de Rytova--Keldysh, dans la forme de Berkelbach, Hybertsen et Reichman~\cite{ch2bhr} :
\[V(r)=\frac{\pi}{2r_0}\,\bigl[\mathbf H_0-Y_0\bigr](r/r_0),\qquad \widehat V(k)=\frac{2\pi}{k(1+r_0k)},\]
où $\mathbf H_0$ est la fonction de Struve, $Y_0$ celle de Neumann, et $r_0>0$ une longueur d'écrantage
(en unités $e^2/(4\pi\varepsilon_0)=1$). Pour $r_0\to0$ on retrouve Coulomb, $2\pi/k$. La forme est citée \tierL{}
(\cite{ch2bhr}, d'après \cite{ch2dlmf} pour $\mathbf K_0=\mathbf H_0-Y_0$) ; la transformée de Fourier est \tierC.
\end{physbox}

Mathlib, à la version épinglée, ne connaît ni Struve ni Neumann. L'idée est donc de \emph{définir} en Lean le potentiel
par une représentation intégrale manifestement complètement monotone --- un mélange de gaussiennes en $r$, c'est-à-dire
un mélange d'exponentielles en $t=r^2$ --- puis de démontrer \emph{à la main}, au niveau C, que cette intégrale est bien
la forme de Struve. L'écart entre les deux est précisément le « trou de fidélité » de cette formalisation.

\paragraph{Étape 0 : le poids.} Pour $x\ge0$ on pose
\[\mathrm{erfcx}(x)=\frac{2}{\sqrt\pi}\int_0^\infty e^{-(y^2+2xy)}\,dy .\]
Le changement de variable $s=y+x$ donne $\mathrm{erfcx}(x)=e^{x^2}\frac2{\sqrt\pi}\int_x^\infty e^{-s^2}ds=e^{x^2}\operatorname{erfc}(x)$,
la fonction d'erreur complémentaire « mise à l'échelle ». Pour $r_0>0$, le poids est $w(a)=r_0^{-1}\,\mathrm{erfcx}(\sqrt a/r_0)$.
\begin{leanbox}[--- \lean{formal_cert/Autoresearch/KeldyshCM.lean:190-191, 244, 248-249}]
\begin{lstlisting}[style=chdeux]
noncomputable def erfcx (x : ℝ) : ℝ :=
  2 / Real.sqrt Real.pi * ∫ y in Ioi (0 : ℝ), Real.exp (-(y ^ 2 + 2 * x * y))
\end{lstlisting}
\begin{lstlisting}[style=chdeux]
noncomputable def w (r0 a : ℝ) : ℝ := 1 / r0 * erfcx (Real.sqrt a / r0)
\end{lstlisting}
\begin{lstlisting}[style=chdeux]
noncomputable def keldysh (r0 t : ℝ) : ℝ :=
  ∫ a in Ioi (0 : ℝ), w r0 a / (2 * a) * Real.exp (-t / (4 * a))
\end{lstlisting}
\end{leanbox}
\noindent Autrement dit, l'objet Lean est
\begin{equation}\label{ch2:eq:kmix}
\mathcal K_{r_0}(t)=\int_0^\infty\frac{w(a)}{2a}\,e^{-t/(4a)}\,da .
\end{equation}

\begin{lemma}[bornes sur le poids]\label{ch2:lem:w}
Pour $x\ge0$, $0<\mathrm{erfcx}(x)\le1$ ; pour $x>0$, $\mathrm{erfcx}(x)\le\frac1{\sqrt\pi\,x}$. Donc, pour $a>0$,
$0<w(a)\le\frac1{r_0}$ et $w(a)\le\frac1{\sqrt\pi\sqrt a}$.
\end{lemma}
\begin{humanproof}
L'intégrande est continu et strictement positif, donc $\mathrm{erfcx}(x)>0$. Pour $y>0$ et $x\ge0$,
$y^2+2xy\ge y^2$, donc l'intégrale est au plus $\int_0^\infty e^{-y^2}dy=\frac{\sqrt\pi}2$, d'où $\mathrm{erfcx}\le1$. Si
$x>0$, on minore plutôt par $2xy$ : $\int_0^\infty e^{-2xy}dy=\frac1{2x}$, d'où
$\mathrm{erfcx}(x)\le\frac2{\sqrt\pi}\cdot\frac1{2x}$. Pour $w$, on substitue $x=\sqrt a/r_0$.
\end{humanproof}

\paragraph{Étape 1 : un lemme générique de mélange de Laplace.} Soit $(\alpha,\mu)$ un espace mesuré, $k,m$ mesurables,
$k\ge0$ et $m\ge0$ presque partout, et supposons $\int k^ne^{-ck}m\,d\mu<\infty$ pour tout $n$ et tout $c>0$. On pose
$G_n(t)=\int(-k)^ne^{-tk}m\,d\mu$.
\begin{theorem}[mélange de Laplace]\label{ch2:thm:laplace}
$g=G_0$ est admissible, et $g^{(n)}(t)=G_n(t)$ pour $t>0$. \tierA
\end{theorem}
\begin{humanproof}
Fixons $t>0$ et faisons varier $s$ dans $]t/2,+\infty[$, voisinage de $t$. La dérivée en $s$ de $(-k)^ne^{-sk}m$ est
$(-k)^{n+1}e^{-sk}m$, dont la valeur absolue vérifie $k^{n+1}e^{-sk}m\le k^{n+1}e^{-(t/2)k}m$ car $s\ge t/2$ et
$k\ge0$. Le majorant ne dépend pas de $s$ et il est intégrable (hypothèse avec $c=t/2$). Le théorème de dérivation sous
le signe intégrale (convergence dominée) donne $G_n'(t)=G_{n+1}(t)$. Par récurrence, $G_0$ est de classe $C^N$ pour tout
$N$ et $g^{(n)}=G_n$. Enfin $(-1)^nG_n(t)=\int k^ne^{-tk}m\,d\mu\ge0$, car l'intégrande est positif presque partout.
\end{humanproof}
\begin{leanbox}[--- \lean{KeldyshCM.lean:31-36, 134-135} (preuve omise)]
\begin{lstlisting}[style=chdeux]
structure MixtureData (μ : Measure α) (k m : α → ℝ) : Prop where
  k_meas : Measurable k
  m_meas : Measurable m
  k_nonneg : ∀ᵐ x ∂μ, 0 ≤ k x
  m_nonneg : ∀ᵐ x ∂μ, 0 ≤ m x
  integrable : ∀ (n : ℕ) (c : ℝ), 0 < c → Integrable (fun x => k x ^ n * Real.exp (-(c * k x)) * m x) μ
\end{lstlisting}
\begin{lstlisting}[style=chdeux]
theorem laplace_admissible (h : MixtureData μ k m) :
    AdmissiblePotential (fun t => ∫ x, Real.exp (-(t * k x)) * m x ∂μ) := by
\end{lstlisting}
\end{leanbox}

Ce théorème est ré-exporté sous le nom \lean{ALEXIS} (chapitre~5).

\paragraph{Étape 2 : l'instance de Keldysh.} On prend la mesure de Lebesgue sur $]0,+\infty[$, $k(a)=\frac1{4a}$ et
$m(a)=\frac{w(a)}{2a}$. Il faut vérifier les moments : $F_n(a)=(4a)^{-n}e^{-c/(4a)}\frac{w(a)}{2a}$ est intégrable.
On coupe en $a=1$. Sur $]0,1]$, avec $u=\frac1{4a}$ : $F_n=2w(a)\,u^{n+1}e^{-cu}$, et $e^{cu}\ge(cu)^N/N!$ donne
$u^Ne^{-cu}\le N!/c^N$ ; $F_n$ est bornée par $2r_0^{-1}(n+1)!/c^{n+1}$ sur un intervalle de longueur finie. Sur
$]1,+\infty[$, $u\le\frac14$, donc $F_n\le\frac{w(a)}{2a}\le\frac1{2\sqrt\pi}a^{-3/2}$ (lemme~\ref{ch2:lem:w}), et
$\int_1^\infty a^{-3/2}da=2$. Le théorème~\ref{ch2:thm:laplace} s'applique.
\begin{leanbox}[--- \lean{KeldyshCM.lean:354} (preuve omise)]
\begin{lstlisting}[style=chdeux]
theorem keldysh_admissible {r0 : ℝ} (hr0 : 0 < r0) : AdmissiblePotential (keldysh r0) := by
\end{lstlisting}
\end{leanbox}

\paragraph{Étape 3 : l'énoncé n'est pas celui de la fonction nulle \tierA.} En Lean, l'intégrale de Bochner d'une
fonction non intégrable vaut $0$, et $0$ est admissible : l'admissibilité seule pourrait être vide. Mais pour $t>0$
l'intégrande de~\eqref{ch2:eq:kmix} est strictement positif et intégrable, donc $\mathcal K_{r_0}(t)>0$.
\begin{leanbox}[--- \lean{KeldyshCM.lean:379} (preuve omise)]
\begin{lstlisting}[style=chdeux]
theorem keldysh_pos {r0 t : ℝ} (hr0 : 0 < r0) (ht : 0 < t) : 0 < keldysh r0 t := by
\end{lstlisting}
\end{leanbox}

\paragraph{Étape 4 : corollaires, et une mise en garde \tierA{} / \tierC.} Le produit par $e^{-\varepsilon t}$
($\varepsilon\ge0$) reste admissible (lemme~\ref{ch2:lem:leibniz}), et le théorème amont donne les deux inégalités
\lean{triangular_keldysh_optimal} et \lean{triangular_keldyshReg_optimal}.
\begin{leanbox}[--- \lean{KeldyshCM.lean:390-392} (énoncé et preuve complète)]
\begin{lstlisting}[style=chdeux]
theorem triangular_keldysh_optimal {r0 : ℝ} (hr0 : 0 < r0) (C : Set Plane) (hC : LocallyFinite C)
    (hdens : DensityOne C) : latticeEnergy (keldysh r0) ≤ energy (keldysh r0) C :=
  (universal_energy_minimum (keldysh r0) C (keldysh_admissible hr0) hC hdens).1
\end{lstlisting}
\end{leanbox}
\begin{warnbox}[: $\top\le\top$]
$V(r)\sim1/r$ à grande distance, et $\sum_{0<|x|\le R}1/|x|$ croît comme $R$ sur un réseau plan. On \emph{s'attend}
donc à $\mathrm{latticeEnergy}(\mathcal K_{r_0})=+\infty$ : la première inégalité se lirait alors $\top\le\top$
\tierC. Ce n'est ni démontré ni réfuté en Lean. Pour $\varepsilon>0$, la somme régularisée devrait être finie \tierC,
mais \textbf{cette finitude n'est pas démontrée}. L'énergie physiquement pertinente d'un cristal de Wigner, neutralisée
par un fond uniforme (« jellium »), n'est couverte par aucun des deux énoncés~\cite{ch2kel}.
\end{warnbox}

\paragraph{Étape 5 : le lien avec la forme de Struve \tierC.}
\begin{theorem}[non vérifié par machine]\label{ch2:thm:struve}
Pour $r_0>0$ et $r>0$ : $V(r)=\mathcal K_{r_0}(r^2)$.
\end{theorem}
\begin{humanproof}
Tous les intégrandes sont positifs : chaque échange d'intégrales est justifié par le théorème de Tonelli.

\emph{(a) Une forme de Laplace en $r$.} La formule DLMF 11.5.2 avec $\nu=0$ donne
$\mathbf K_0(z)=\frac2\pi\int_0^\infty e^{-zu}(1+u^2)^{-1/2}du$ pour $\operatorname{Re}z>0$ \tierL~\cite{ch2dlmf}. Avec
$z=r/r_0$ : $V(r)=\frac1{r_0}\int_0^\infty e^{-ru/r_0}(1+u^2)^{-1/2}du$. Cela suffit déjà à voir que $V$ est CM \emph{en $r$}.

\emph{(b) Subordination.} Pour $\rho>0$ et $t>0$,
$e^{-\rho\sqrt t}=\int_0^\infty e^{-ta}\,\frac{\rho}{2\sqrt\pi}\,a^{-3/2}e^{-\rho^2/(4a)}\,da$ (transformée de Laplace de
la densité de Lévy, formule classique utilisée sans citation dans l'article source). On prend $\rho=u/r_0$, $t=r^2$, on
insère dans (a) et l'on échange les intégrales :
\[V(r)=\int_0^\infty e^{-ta}\,\frac{a^{-3/2}}{2\sqrt\pi\,r_0^2}\left[\int_0^\infty\frac{u}{\sqrt{1+u^2}}e^{-\lambda u^2}du\right]da,
\qquad\lambda=\frac1{4r_0^2a}.\]

\emph{(c) L'intégrale intérieure est élémentaire.} Avec $v=\sqrt{1+u^2}$, $dv=u(1+u^2)^{-1/2}du$ :
\[\int_0^\infty\frac{u}{\sqrt{1+u^2}}e^{-\lambda u^2}du=e^{\lambda}\int_1^\infty e^{-\lambda v^2}dv
=\frac12\sqrt{\frac\pi\lambda}\,\mathrm{erfcx}(\sqrt\lambda).\]
Avec $\sqrt{\pi/\lambda}=2r_0\sqrt{\pi a}$, on obtient $V(r)=\int_0^\infty e^{-ta}\rho(a)\,da$ où
$\rho(a)=\dfrac{\mathrm{erfcx}\bigl(1/(2r_0\sqrt a)\bigr)}{2r_0\,a}$.

\emph{(d) Changement de variable.} On pose $a=\frac1{4b}$, d'où $|da|=\frac{db}{4b^2}$ et $e^{-ta}=e^{-t/(4b)}$. Alors
$\rho(a)\,da=\frac{4b\,\mathrm{erfcx}(\sqrt b/r_0)}{2r_0}\cdot\frac{db}{4b^2}=\frac{w(b)}{2b}\,db$, donc
$V(r)=\int_0^\infty\frac{w(b)}{2b}e^{-r^2/(4b)}db=\mathcal K_{r_0}(r^2)$.
\end{humanproof}
\noindent Des vérifications numériques à 30 chiffres (mpmath) confirment chaque étape, avec deux erreurs plantées
détectées \tierX~\cite{ch2kel} ; elles ne remplacent pas la démonstration. Le lecteur doit accepter le
théorème~\ref{ch2:thm:struve} pour transférer l'énoncé Lean, qui porte sur~\eqref{ch2:eq:kmix}, à la forme de Struve.
Article : doi\string:10.5281/zenodo.23646107.

% ---------------------------------------------------------------------
\section{La généralisation : une bibliothèque CM/Bernstein}\label{ch2:sec:lib}

Les trois démonstrations précédentes ont un air de famille : une différence de translatées (dipôles), un produit et une
composition avec $\sqrt t$ (Yukawa), un mélange de Laplace (Keldysh). La bibliothèque P3~\cite{p3} (doi\string:10.5281/zenodo.23653842)
isole ces mécanismes une fois pour toutes. Bilan de la porte : \textbf{130 théorèmes} déclarés, chacun avec sa ligne
\texttt{\#print axioms}, tous dans la liste autorisée, aucun avertissement ; un audit par un agent d'IA séparé a
reconstruit les 11 fichiers et 5 modules prérequis depuis les sources et trouvé 0 constante hors liste parmi 167
\tierA. Aucun audit humain ; les 20 drapeaux \texttt{LEDGER\_UNAUDITED\_TIER\_A} restent ouverts.

\subsection{Deux définitions et un pont}
\begin{definition}[fonction de Bernstein]
$\psi:\,]0,+\infty[\,\to\R$ est \emph{de Bernstein} si elle est $C^\infty$, positive au sens large ($\psi\ge0$), et si $\psi'$ est CM.
Exemples : $at+b$ ($a,b\ge0$), $t^\alpha$ pour $0\le\alpha\le1$, en particulier $\sqrt t$ ; $1-e^{-at}$ ; $\ln(1+t)$.
\end{definition}
\begin{leanbox}[--- \lean{formal_cert/Programmes/CompletelyMonotone/Basic.lean:28-30, 33-34} et \lean{Bernstein.lean:24-25}]
\begin{lstlisting}[style=chdeux]
def IsCompletelyMonotoneOn (f : ℝ → ℝ) : Prop :=
  ContDiffOn ℝ (⊤ : ℕ∞) f (Set.Ioi 0) ∧
  ∀ (n : ℕ) (t : ℝ), 0 < t → 0 ≤ (-1 : ℝ) ^ n * iteratedDeriv n f t
\end{lstlisting}
\begin{lstlisting}[style=chdeux]
theorem isCompletelyMonotoneOn_iff_admissible {f : ℝ → ℝ} :
    IsCompletelyMonotoneOn f ↔ OAI.AtomicTriangular.AdmissiblePotential f := by
\end{lstlisting}
\begin{lstlisting}[style=chdeux]
def IsBernsteinOn (g : ℝ → ℝ) : Prop :=
  ContDiffOn ℝ (⊤ : ℕ∞) g (Set.Ioi 0) ∧ (∀ t : ℝ, 0 < t → 0 ≤ g t) ∧ IsCompletelyMonotoneOn (deriv g)
\end{lstlisting}
\end{leanbox}
\noindent Le pont \lean{isCompletelyMonotoneOn_iff_admissible} est une \emph{équivalence} \tierA : la clause $g\ge0$ de
l'amont est le cas $n=0$ de la troisième. Toute fonction CM de la bibliothèque est donc directement utilisable dans
\lean{universal_energy_minimum}.

\subsection{Les lemmes de clôture}
On note $h_n(t)=(-1)^nf^{(n)}(t)$ pour $f$ CM.
\begin{enumerate}[label=(\roman*),itemsep=2pt]
\item \emph{Cône.} Les dérivées sont linéaires : $cf$ ($c\ge0$), $f+g$, et toute combinaison finie positive de
fonctions CM est CM. Les constantes positives sont CM.
\item \emph{Produit.} Par le lemme de Leibniz~\ref{ch2:lem:leibniz}, $fg$ est CM.
\item \emph{Dérivée.} $h_n'=(-1)^nf^{(n+1)}=-h_{n+1}\le0$, donc chaque $h_n$ décroît (accroissements finis). En
particulier $-f'$ est CM, car $(-1)^n(-f')^{(n)}=h_{n+1}\ge0$.
\item \emph{Translatée.} Pour $c\ge0$, $t\mapsto f(t+c)$ est CM : ses dérivées sont $f^{(n)}(t+c)$ et $t+c>0$.
\item \emph{Différence de translatées.} Pour $c\ge0$, $f(t)-f(t+c)$ est CM, car
$(-1)^n\bigl[f(t)-f(t+c)\bigr]^{(n)}=h_n(t)-h_n(t+c)\ge0$ puisque $h_n$ décroît. Aucune représentation intégrale n'est utilisée.
\end{enumerate}
\begin{leanbox}[--- \lean{Basic.lean:160-161} (preuve omise ; elle tient en dix lignes)]
\begin{lstlisting}[style=chdeux]
theorem sub_comp_add_const (hf : IsCompletelyMonotoneOn f) {c : ℝ} (hc : 0 ≤ c) :
    IsCompletelyMonotoneOn (fun t => f t - f (t + c)) := by
\end{lstlisting}
\end{leanbox}
\noindent Le point (v) est la version abstraite de la démonstration dipolaire : le calcul de Pochhammer y est remplacé par
la seule décroissance de $h_n$.

\subsection{Le théorème de subordination}
\begin{theorem}[subordination]\label{ch2:thm:subord}
Si $f$ est CM et si $\psi$ est une fonction de Bernstein strictement positive sur $]0,+\infty[$, alors $f\circ\psi$ est
CM. \tierA{} ; démonstration ci-dessous \tierC.
\end{theorem}
\begin{humanproof}
La régularité vient de la stabilité de $C^\infty$ par composition ($\psi$ envoie $]0,+\infty[$ dans $]0,+\infty[$).
Pour les signes, on démontre par récurrence sur $n$ la propriété
\[\mathcal Q(n)\ :\ \text{« pour \emph{toute} fonction CM } f\text{ et tout } m\le n,\ (-1)^m(f\circ\psi)^{(m)}\ge0\text{ sur }]0,+\infty[\text{ ».}\]
$\mathcal Q(0)$ : $f(\psi(t))\ge0$ car $\psi(t)>0$ et $f\ge0$.

$\mathcal Q(n)\Rightarrow\mathcal Q(n+1)$ : soit $f$ CM ; les ordres $m\le n$ sont couverts par $\mathcal Q(n)$. Pour
$m=n+1$, posons $F=-f'$, qui est CM d'après (iii). La dérivation composée donne
$(f\circ\psi)'=(f'\circ\psi)\,\psi'=-(F\circ\psi)\,\psi'$, donc
\[(-1)^{n+1}(f\circ\psi)^{(n+1)}=(-1)^n\bigl[(F\circ\psi)\,\psi'\bigr]^{(n)}\ge0\]
par le lemme de Leibniz, appliqué à $u=F\circ\psi$ (dont les signes jusqu'à l'ordre $n$ sont donnés par $\mathcal Q(n)$
\emph{appliquée à $F$}) et à $v=\psi'$ (CM par définition d'une fonction de Bernstein). C'est parce que $\mathcal Q(n)$
quantifie sur \emph{toutes} les fonctions CM que la récurrence se referme sans la formule de Faà di Bruno.
\end{humanproof}
\begin{leanbox}[--- \lean{Bernstein.lean:49-51} (preuve omise) et \lean{:85-87} (énoncé et preuve complète)]
\begin{lstlisting}[style=chdeux]
theorem comp_sign (hg : IsBernsteinOn g) (hpos : ∀ t, 0 < t → 0 < g t) :
    ∀ n : ℕ, ∀ f : ℝ → ℝ, IsCompletelyMonotoneOn f → ∀ m ≤ n, ∀ t : ℝ, 0 < t →
      0 ≤ (-1 : ℝ) ^ m * iteratedDeriv m (fun t => f (g t)) t := by
\end{lstlisting}
\begin{lstlisting}[style=chdeux]
theorem IsCompletelyMonotoneOn.comp_bernstein {f g : ℝ → ℝ} (hf : IsCompletelyMonotoneOn f)
    (hg : IsBernsteinOn g) (hpos : ∀ t, 0 < t → 0 < g t) : IsCompletelyMonotoneOn (fun t => f (g t)) :=
  ⟨hf.1.comp hg.1 (fun s hs => hpos s hs), fun n t ht => hg.comp_sign hpos n f hf n le_rfl t ht⟩
\end{lstlisting}
\end{leanbox}
\noindent L'énoncé Lean \lean{comp_sign} porte « $\forall f$ » \emph{à l'intérieur} de la récurrence sur $n$ : c'est
littéralement la propriété $\mathcal Q(n)$.

\begin{remark}[l'hypothèse $\psi>0$ est nécessaire \tierA]
$f$ n'est contrainte que sur $]0,+\infty[$, et $\psi\equiv0$ est de Bernstein ; alors $f\circ\psi$ est la constante
« poubelle » $f(0)$, qui peut être négative. Le témoin \lean{junk_comp_zero_not_cm} réalise ce cas. Le théorème est
classique dans l'esprit (« subordination de Bochner ») ; l'article source n'a pas localisé son numéro dans la
monographie de référence~\cite{ch2ssv}. La contribution est la preuve vérifiée par machine sur $]0,+\infty[$.
\end{remark}

\paragraph{Instances.} $\psi(t)=t^\alpha$ ($0\le\alpha\le1$) est de Bernstein : $\psi'=\alpha t^{-(1-\alpha)}$ est CM
par le lemme~\ref{ch2:lem:riesz}. D'où, avec $f(u)=e^{-au}$ : $e^{-at^\alpha}$ est CM (exponentielle étirée), et avec
$\alpha=\frac12$ : $e^{-\kappa\sqrt t}$ est CM pour tout $\kappa\ge0$, en une ligne.
\begin{leanbox}[--- \lean{Bernstein.lean:160-161} (énoncé et preuve complète)]
\begin{lstlisting}[style=chdeux]
theorem phi_cm {κ : ℝ} (hκ : 0 ≤ κ) : IsCompletelyMonotoneOn (AutoresearchYukawa.phi κ) :=
  ((exp_neg_mul hκ).comp_sqrt).congr' fun t _ => by simp only [AutoresearchYukawa.phi, neg_mul]
\end{lstlisting}
\end{leanbox}
\noindent Comparez avec la récurrence forte de la proposition~\ref{ch2:prop:phi} : celle-ci exigeait $\kappa>0$ et une
induction ad hoc ; la bibliothèque couvre $\kappa\ge0$.

\paragraph{Mélanges de Laplace, hypothèse affaiblie.} Le théorème \lean{laplace} de \lean{Laplace.lean:49} ne suppose
plus que tous les moments sont finis, mais seulement que $e^{-ck}m$ est intégrable pour tout $c>0$ : l'inégalité
$u^n\le n!\,(c/2)^{-n}e^{(c/2)u}$ ($u\ge0$) donne $k^ne^{-ck}m\le\frac{n!}{(c/2)^n}e^{-(c/2)k}m$, d'où tous les
moments. C'est la direction \emph{facile} du théorème de Bernstein--Widder (une transformée de Laplace d'une mesure
positive est CM) ; la direction difficile (toute fonction CM est une telle transformée~\cite{widder}) n'est
\emph{pas} démontrée.

\subsection{Les trois potentiels deviennent des corollaires}
\begin{leanbox}[--- \lean{Corollaries.lean:24-25, 28-30, 38}]
\begin{lstlisting}[style=chdeux]
theorem yukawa_cm {κ : ℝ} (hκ : 0 ≤ κ) : IsCompletelyMonotoneOn (AutoresearchYukawa.g κ) :=
  ((riesz_cm zero_le_one).mul (phi_cm hκ)).congr (AutoresearchYukawa.g_eqOn κ)
\end{lstlisting}
\begin{lstlisting}[style=chdeux]
theorem dipolar_cm (d : ℝ) : IsCompletelyMonotoneOn (AutoresearchDipolar.g d) :=
  (((rpow_neg (by norm_num : (0 : ℝ) ≤ 1 / 2)).sub_comp_add_const (sq_nonneg d)).const_mul
    zero_le_two).congr' fun _ _ => rfl
\end{lstlisting}
\begin{lstlisting}[style=chdeux]
theorem keldysh_cm {r0 : ℝ} (hr0 : 0 < r0) : IsCompletelyMonotoneOn (AutoresearchKeldysh.keldysh r0) := by
\end{lstlisting}
\end{leanbox}
\noindent Lecture : Yukawa $=$ Riesz $\times$ ($\exp\circ\sqrt{\ }$) ; dipôle $=2\times$ (différence de translatées de
$t^{-1/2}$) ; Keldysh $=$ mélange de Laplace (preuve de huit lignes, omise).

\begin{warnbox}[: ce que la bibliothèque ne raccourcit pas]
L'expression « corollaires en une ligne » n'est vraie qu'à moitié. L'article P3 a \emph{mesuré} la fermeture de
dépendances des trois résultats H1, H2, H3 réunis : elle \emph{augmente} quand on passe par la bibliothèque. Seul le
coût \emph{marginal} d'un nouveau potentiel devient petit. Pour Keldysh, la preuve courte s'appuie toujours sur les
bornes lourdes de \lean{keldysh_data} (erfcx, intégrabilité) du module d'origine : l'hypothèse affaiblie est
\emph{disponible} mais n'est pas \emph{exercée}. Mathlib contient déjà une notion voisine, \lean{AbsolutelyMonotoneOn}
(sans signes alternés) ; la bibliothèque n'est pas prête pour Mathlib~\cite{p3}.
\end{warnbox}

\paragraph{Nouvelles conséquences \tierA, conditionnelles pour l'optimalité.} Yukawa à c\oe{}ur mou
$g_\kappa(t+b^2)$ (translatée), exponentielle étirée $e^{-at^\alpha}$, exciton dipolaire écranté
$2(g_\kappa(t)-g_\kappa(t+d^2))$, mélanges continus de Yukawa. Un exemple instructif, les quadrupôles $+1,-2,+1$ dans
les couches $0,d,2d$ ($c=d^2$) :
\begin{itemize}
\item le \emph{champ} d'un quadrupôle sur une charge test, $f(t)-2f(t+c)+f(t+4c)$ avec $f=t^{-1/2}$, n'est \emph{pas}
CM : en $t=3$, $c=1$, il vaut $\frac1{\sqrt3}-1+\frac1{\sqrt7}<0$ (\lean{quadrupole_not_cm}) ;
\item l'\emph{interaction} de deux quadrupôles alignés compte les paires ($1+4+1=6$ à $\Delta z=0$, $4\times(-2)=-8$ à
$|\Delta z|=d$, $2\times1=2$ à $|\Delta z|=2d$) et vaut $2\bigl(3f(t)-4f(t+c)+f(t+4c)\bigr)$ \tierC{} pour le comptage ;
elle \emph{est} CM pour toute $f$ CM (exercice~\ref{ch2:ex:quad}).
\end{itemize}
\begin{leanbox}[--- \lean{Applications.lean:172, 175-176} (preuve omise)]
\begin{lstlisting}[style=chdeux]
noncomputable def quadrupolarPair (f : ℝ → ℝ) (c t : ℝ) : ℝ := 2 * (3 * f t - 4 * f (t + c) + f (t + 4 * c))
\end{lstlisting}
\begin{lstlisting}[style=chdeux]
theorem quadrupolarPair_cm {f : ℝ → ℝ} (hf : IsCompletelyMonotoneOn f) {c : ℝ} (hc : 0 ≤ c) :
    IsCompletelyMonotoneOn (quadrupolarPair f c) := by
\end{lstlisting}
\end{leanbox}
\noindent Pour ces nouveaux potentiels, aucune finitude de l'énergie optimale n'est démontrée : pour des queues lentes,
l'inégalité d'optimalité peut se lire $\top\le\top$.

% ---------------------------------------------------------------------
\section{Un problème voisin : la conjecture de Sandier--Serfaty}\label{ch2:sec:ss}

\begin{warnbox}[: statut]
\textbf{La conjecture de Sandier--Serfaty n'est pas démontrée dans ce livre.} Un manuscrit de la publication openai/math
en revendique une preuve ; ce manuscrit est produit par une IA et n'est pas relu par des pairs. Ce que nous avons
fait --- reproduire un certificat, auditer une réduction, écrire un énoncé Lean, formaliser une petite partie --- ne
constitue ni une preuve, ni une vérification, ni une relecture par des pairs.
\end{warnbox}

\subsection{Ce que dit la conjecture}\label{ch2:ssec:ssconj}
\begin{physbox}[: des vortex aux réseaux \tierC]
Dans un supraconducteur de type II soumis à un champ magnétique, le champ pénètre sous forme de tubes de flux
quantifiés, les \emph{vortex}, qui s'organisent en réseau (réseau d'Abrikosov, observé triangulaire). Sandier et Serfaty
ont déduit du modèle de Ginzburg--Landau une \emph{énergie renormalisée} $W$ qui gouverne la position des vortex à
l'ordre suivant~\cite{ch2ss} \tierL. Le lien avec la physique des vortex est ici une glose \tierC.
\end{physbox}
Voici la formulation du manuscrit, qui suit Sandier--Serfaty \tierL. Soit $\Lambda\subset\R^2$ localement fini, de
densité bornée, et $\nu_\Lambda=\sum_{p\in\Lambda}\delta_p$. Un champ $E:\R^2\to\R^2$ localement intégrable est
\emph{admissible} ($E\in\mathcal A_1$) si, au sens des distributions,
\[\operatorname{div}E=2\pi(\nu_\Lambda-1),\qquad\operatorname{rot}E=0 .\]
Chaque point de $\Lambda$ est une charge unité, le « $-1$ » est un fond neutralisant uniforme. Près d'une charge,
$|E|^2\sim1/|x-p|^2$ n'est pas intégrable : on retire des disques de rayon $\eta$ et l'on compense,
\[W(E,\chi_R)=\lim_{\eta\to0}\Bigl[\tfrac12\int_{\R^2\setminus\bigcup_pB(p,\eta)}\chi_R|E|^2+\pi\ln\eta\sum_{p\in\Lambda}\chi_R(p)\Bigr],
\qquad W(E)=\limsup_{R\to\infty}\frac{W(E,\chi_R)}{4R^2},\]
où $\chi_R$ est une troncature lisse à support dans le carré $[-R,R]^2$. Le signe $+\pi\ln\eta$ est le bon :
$\frac12\int_{\eta<|x-p|<\rho}|x-p|^{-2}dx=\pi\ln\rho-\pi\ln\eta$.
\begin{itemize}
\item \textbf{Conjecture} (Sandier--Serfaty 2012, Conjecture~1) : $W(E)\ge W(E_\triangle)$ pour tout $E\in\mathcal A_1$,
où $E_\triangle$ est le champ périodique du réseau triangulaire $\Lambda_\triangle$ de covolume un \tierL.
\item Sandier--Serfaty la démontrent \emph{parmi les réseaux} (leur Théorème~2) et écrivent que la minimisation en
général reste ouverte \tierL~\cite{ch2ss}.
\item Bétermin--Sandier montrent qu'elle équivaut à une valeur conjecturée par Brauchart--Hardin--Saff pour le terme
d'ordre suivant de l'énergie logarithmique optimale de $N$ points sur la sphère~\cite{ch2bs,ch2bhs} \tierL, et donnent
$W(E_\triangle)=\pi\ln\frac{2\sqrt2\,\pi}{\sqrt3\,\Gamma(1/3)^3}\approx-4{,}1504128$ \tierL.
\end{itemize}

\paragraph{Ce que revendique le manuscrit openai/math.} Une preuve~\cite{ch2oaiss}, par (1) une réduction du plan aux
tores carrés $\R^2/(\sqrt n\,\Z^2)$, (2) un potentiel d'essai recollé sur les secteurs de Voronoï, \emph{exactement
calibré} au réseau triangulaire, et (3) un certificat fini en arithmétique d'intervalles exacte (entiers de 144~bits),
organisé en douze étapes. Ce manuscrit est produit par une IA et n'est pas relu par des pairs.

\subsection{P1 : reproduire et auditer, sans démontrer}\label{ch2:ssec:p1}
\begin{itemize}
\item \textbf{Reproduction du certificat \tierX.} Le programme C++ inchangé a été recompilé et exécuté sur une autre
chaîne d'outils et une autre machine (g++~11.4) ; l'exécution du certificat a pris environ 478~s ; la ligne finale est
\texttt{OK} et les 22 paires d'entiers imprimées sont identiques, 22 sur 22, à l'enregistrement d'origine.
\emph{Témoins négatifs} : trois versions affaiblies d'une ligne ; deux sont rejetées par les assertions du programme,
\textbf{la troisième passe} (les issues des deux dernières avaient été prédites par écrit avant la fin des calculs). La
reproduction montre que le programme se comporte comme enregistré ; elle ne teste aucune des bornes analytiques qu'il
est censé certifier.
\item \textbf{Audit de la réduction \tierC.} Un agent d'IA a vérifié pas à pas les sections 2 à 5 du manuscrit
(31 étapes numérotées et 14 lignes de l'étape certificat) : \emph{aucune étape n'a été classée comme lacune}, avec des
limites déclarées (les morceaux de spline hors de $[p,R_0]$, les valeurs de l'étape de grille et la confiance dans
l'exécution ne sont pas vérifiés indépendamment). Ce n'est ni une preuve ni une relecture par des pairs.
\item \textbf{Énoncé Lean verrouillé.} La conjecture est écrite comme une proposition Lean (\lean{Prop}),
\textbf{non démontrée}, dont la région d'énoncé est verrouillée par empreinte (\texttt{statement\_lock.py}).
\end{itemize}
\begin{leanbox}[--- \lean{formal_cert/Programmes/SandierSerfatyStatement.lean:153-160} (une définition, pas un théorème)]
\begin{lstlisting}[style=chdeux]
def SandierSerfatyStatement : Prop :=
  (∃ χ, IsSquareCutoffFamily χ) ∧
    ∃ E₀ : Plane → Plane, Admissible triLattice E₀ ∧ IsPeriodic triLattice E₀ ∧
      ∀ χ : ℝ → Plane → ℝ, IsSquareCutoffFamily χ →
        ∃ w₀ : ℝ → ℝ, (∀ R : ℝ, 1 < R → HasLocalEnergy triLattice E₀ (χ R) (w₀ R)) ∧
          WSq w₀ ≠ ⊤ ∧ WSq w₀ ≠ ⊥ ∧
          ∀ (Λ : Set Plane) (E : Plane → Plane) (w : ℝ → ℝ), Admissible Λ E →
            (∀ R : ℝ, 1 < R → HasLocalEnergy Λ E (χ R) (w R)) → WSq w₀ ≤ WSq w
\end{lstlisting}
\end{leanbox}
\noindent Seuls dix lemmes de non-dégénérescence sont \tierAs{} : la classe admissible est non vide (le champ $E=-\pi x$,
sans charge), la constante de fond $2\pi$ est porteuse (avec $0$ ou $-2\pi$ à la place, ce champ n'est plus
admissible), les classes de troncatures sont non vides. Attention : le témoin de \lean{competitor_nonvacuous} est ce
champ $-\pi x$, dont l'énergie est $+\infty$ \tierC{} ; il montre que les hypothèses sur un concurrent sont
\emph{satisfiables}, pas que l'inégalité contraint un concurrent. L'existence de la limite en $\eta$ est \emph{supposée}
(\lean{HasLocalEnergy}) et le minimiseur triangulaire est un témoin existentiel : ce sont des écarts déclarés à la
formulation du manuscrit~\cite{ch2ssa}. Article P1 : doi\string:10.5281/zenodo.23653827.

\subsection{A1 : formaliser une partie, et mesurer le reste}\label{ch2:ssec:a1}
\paragraph{Arithmétique d'intervalles en Lean.} Un intervalle est un couple d'entiers $(u,v)$ lu comme $[u/Q,v/Q]$ avec
$Q=2^{144}$. Les produits sont arrondis \emph{vers l'extérieur} : borne inférieure par $\lfloor\cdot/Q\rfloor$, borne
supérieure par $\lceil\cdot/Q\rceil$, de sorte que le vrai résultat réel n'est jamais perdu. Exemple jouet avec $Q=100$ :
$[1{,}20;1{,}30]\times[0{,}50;0{,}60]$, soit $(120,130)\times(50,60)$, a pour produits d'extrémités
$6000,7200,6500,7800$, d'où $(\lfloor6000/100\rfloor,\lceil7800/100\rceil)=(60,78)=[0{,}60;0{,}78]$~\cite{ch2ssl}.

\paragraph{Deux niveaux de vérification} (axe distinct des niveaux A/L/C/X) :
\textbf{L1} (réexécution par le noyau) : la transcription Lean d'une étape, évaluée par le noyau avec
\lean{decide +kernel} (et non \lean{native_decide}, qui ajoute un axiome et fait confiance au code compilé), redonne
bit pour bit les entiers imprimés par le C++. \textbf{L2} (énoncé réel) : un théorème Lean affirme qu'un \emph{nombre
réel} appartient à l'intervalle.

\begin{center}\small
\begin{tabular}{@{}p{0.27\linewidth}p{0.4\linewidth}p{0.29\linewidth}@{}}\toprule
Résultat (journaux sauvegardés) & Valeur & Niveau\\\midrule
Étapes complètes en L1 & 4 sur 12 (\lean{init_constants}, \lean{bcons/eps}, \lean{calcLam}, \lean{calcK}) & \tierAs{} pour \lean{init_constants}, \lean{bcons/eps} ; \tierC{} pour \lean{calcLam}, \lean{calcK} (pas de recompilation indépendante)\\
\quad dont conditionnelles & \lean{calcLam}, \lean{calcK} (tables en prémisses), \lean{bcons} (coefficients $c_m$) & \\
Étapes complètes en L2 & 1 sur 12 (\lean{init_constants} : $\ln2,\pi,p,R_0,a$ réels) & \tierAs\\
Part des opérations d'intervalles réexécutées & $1\,685\,129$ sur $1\,238\,829\,689$, soit $0{,}136\,\%$ & \tierX\\
Énoncés faux plantés rejetés & 11 & \tierX\\
Coût projeté du reste en L1 & 854 à 2\,240 heures-c\oe{}ur & \tierX\\\bottomrule
\end{tabular}
\end{center}
Les 22 paires imprimées et les assertions finales du C++ sont en outre vérifiées comme arithmétique sur les extrémités
exactes (niveau « E », plus faible), sauf le test de borne de l'étape \lean{core_test} ; le test de séparation des
tables n'est vérifié que fenêtre par fenêtre. Vérification indépendante : l'audit par un second agent d'IA a reconstruit
13 modules Lean depuis les sources, mais \emph{pas} \lean{CalcK.lean}, \lean{CalcLam.lean} ni 28 fichiers de morceaux ;
les résultats L1 de \lean{calcLam} et \lean{calcK} n'ont donc pas eu de passe producteur${}\ne{}$vérificateur
\tierC~\cite{ch2ssl}. Le travail restant le plus lourd
est le niveau L2 du contenu analytique (fonctions radiales, coefficients de réseau, restes de quadrature), puis
30 autres étapes de la réduction : un projet de plusieurs mois \tierC~\cite{ch2ssl}. Article A1 :
doi\string:10.5281/zenodo.23656804.

\paragraph{Le lemme \lean{signed_endpoints} \tierAs.} C'est le seul lemme de la réduction démontré complètement. Il
permet de traiter un secteur de Voronoï dont le pied de la perpendiculaire tombe hors de son arête. Cadre : un noyau
$G(d,s)$, pair en $s$ et localement intégrable pour $d\ge d_0$, et sa primitive $g_d(x)=\int_0^xG(d,s)\,ds$.

\begin{center}\small
\begin{tabular}{@{}p{0.34\linewidth}p{0.6\linewidth}@{}}\toprule
Manuscrit & Lean (\lean{SignedEndpoints.lean:79-85}, verbatim)\\\midrule
$K<0$, $d_0>0$ ; $G(d,\cdot)$ pair & \lstinline[style=chdeux]!(hK : K < 0) (hd₀ : 0 < d₀)!, \lstinline[style=chdeux]!heven : ∀ d, d₀ ≤ d → ∀ s, G d (-s) = G d s!\\
localement intégrable & \lstinline[style=chdeux]!hint : ∀ d, d₀ ≤ d → ∀ a b : ℝ, IntervalIntegrable (G d) volume a b!\\
$2g_d(x)\ge K$ pour $x\ge0$ & \lstinline[style=chdeux]!hpos : ∀ d, d₀ ≤ d → ∀ x, 0 ≤ x → K ≤ 2 * gPrim G d x!\\
$\int_0^\infty\max\{0,-G(d,s)\}\,ds\le-K$ & \lstinline[style=chdeux]!hneg : ∀ d, d₀ ≤ d → ∫⁻ s in Ioi 0, ENNReal.ofReal (-G d s) ≤ ENNReal.ofReal (-K)!\\
conclusion pour $|\alpha|,|\beta|<\frac\pi2$, $\alpha+\beta>0$ & \lstinline[style=chdeux]!K ≤ gPrim G d (d * Real.tan α) + gPrim G d (d * Real.tan β)!\\\bottomrule
\end{tabular}
\end{center}
\begin{leanbox}[--- \lean{formal_cert/Programmes/SandierSerfaty/SignedEndpoints.lean:33, 79-85} (preuve omise)]
\begin{lstlisting}[style=chdeux]
noncomputable def gPrim (G : ℝ → ℝ → ℝ) (d x : ℝ) : ℝ := ∫ s in (0 : ℝ)..x, G d s
\end{lstlisting}
\begin{lstlisting}[style=chdeux]
theorem signed_endpoints (G : ℝ → ℝ → ℝ) (K d₀ : ℝ) (hK : K < 0) (hd₀ : 0 < d₀)
    (heven : ∀ d, d₀ ≤ d → ∀ s, G d (-s) = G d s)
    (hint : ∀ d, d₀ ≤ d → ∀ a b : ℝ, IntervalIntegrable (G d) volume a b)
    (hpos : ∀ d, d₀ ≤ d → ∀ x, 0 ≤ x → K ≤ 2 * gPrim G d x)
    (hneg : ∀ d, d₀ ≤ d → ∫⁻ s in Ioi 0, ENNReal.ofReal (-G d s) ≤ ENNReal.ofReal (-K)) :
    ∀ d, d₀ ≤ d → ∀ α β : ℝ, |α| < π / 2 → |β| < π / 2 → 0 < α + β →
      K ≤ gPrim G d (d * Real.tan α) + gPrim G d (d * Real.tan β) := by
\end{lstlisting}
\end{leanbox}
\noindent L'intégrale de la partie négative est une \lean{lintegral} ($\int^-$, à valeurs dans $[0,+\infty]$) : une
intégrale infinie \emph{viole} l'hypothèse, au lieu de produire la valeur poubelle $0$ de l'intégrale de Bochner.

\begin{theorem}[points d'extrémité signés]\label{ch2:thm:signed}
Sous les hypothèses du tableau, pour $d\ge d_0$ et $|\alpha|,|\beta|<\frac\pi2$ avec $\alpha+\beta>0$ :
$g_d(d\tan\alpha)+g_d(d\tan\beta)\ge K$. \tierAs
\end{theorem}
\begin{humanproof}
\emph{(1) Imparité.} Par parité de $G(d,\cdot)$ et le changement de variable $s\mapsto-s$, $g_d(-x)=-g_d(x)$.

\emph{(2) Deux angles positifs.} Alors $x=d\tan\alpha\ge0$ et $y=d\tan\beta\ge0$ ; l'hypothèse $2g_d\ge K$ donne
$g_d(x)\ge K/2$ et $g_d(y)\ge K/2$ ; on additionne.

\emph{(3) Signes opposés, $\alpha<0$.} Alors $\alpha+\beta>0$ force $\beta>-\alpha>0$. Posons $u=-d\tan\alpha>0$ et
$v=d\tan\beta$ ; $\tan$ étant strictement croissante sur $]-\frac\pi2,\frac\pi2[$, on a $u<v$. Par imparité,
$g_d(d\tan\alpha)+g_d(d\tan\beta)=g_d(v)-g_d(u)=\int_u^vG(d,s)\,ds$.

\emph{(4) Estimation centrale.} $\int_u^vG\ge-\int_u^vG^-\ge-\int_0^\infty G^-\ge K$, où $G^-=\max(0,-G)$, par
l'hypothèse sur la partie négative et $0<u$.

\emph{(5)} Le cas $\beta<0$ est symétrique ; les deux négatifs sont exclus par $\alpha+\beta>0$.
\end{humanproof}
\noindent Sans $\alpha+\beta>0$, l'énoncé est faux (exercice~\ref{ch2:ex:signed}) ; la réfutation est vérifiée par le
noyau (\lean{ctl_signed_without_sum_false}), et le même noyau $G\equiv1$ satisfait toutes les autres hypothèses
(\lean{ctl_signed_hyps}), si bien que le théorème n'est pas vide.

\begin{warnbox}[: ce que \lean{signed_endpoints} ne fait pas]
Le lemme est énoncé pour un $G$ \emph{abstrait}. Le $G$ concret demande les données radiales $f_0,e$ de la section~4 du
manuscrit, qui ne sont pas définies en Lean, et rien ne déduit ses deux hypothèses du certificat : le lemme \textbf{n'est
pas relié au certificat}. Le critère d'arrêt préenregistré s'applique : \textbf{aucune preuve formelle de la conjecture
n'est revendiquée}. La fidélité de l'énoncé au manuscrit a été comparée par un agent, pas par un humain.
\end{warnbox}

\subsection{B3 : pourquoi la méthode ne passe pas telle quelle en dimension trois}\label{ch2:ssec:b3}
En dimension trois, l'analogue est la conjecture de Wigner : le réseau cubique centré (bcc) minimiserait l'énergie du
jellium parmi \emph{toutes} les configurations ; il n'est connu optimal que parmi les réseaux de Bravais de covolume un, pour l'énergie de Born coulombienne (Muratov--Novaga, théorème~1, p.~4~\cite{ch2mn})
\tierL. L'étude B3 (doi\string:10.5281/zenodo.23655002)~\cite{ch2b3} découpe la preuve plane en 13 étapes et classe chacune :
indépendante de la dimension, transférable, ou spécifiquement plane. Le décompte d'incidences (identité d'Euler par
cellule) \emph{n'est pas} l'obstacle.

\paragraph{L'obstacle : la calibration exacte.} Toute la preuve plane est une minoration par un potentiel d'essai $H$ :
$W(h)=F(H)+\frac1{8\pi}\int|\nabla(h-H)|^2\ge F(H)$ (identité duale, \tierC), et la minoration n'est optimale que si
$H$ \emph{égale} la vraie fonction de Green $h$ au réseau optimal. En dimension deux c'est possible : dans chaque secteur,
$\cos6m\varphi=T_{6m}(d/t)$ (polynôme de Tchebychev, $t=d/\cos\varphi$), donc $h_\triangle$ ne dépend que de deux
variables et appartient à la classe d'essai. En dimension trois, la restriction de la fonction de Green bcc à un cercle
de rayon fixé sur une face porte un mode $\cos6\psi$ (faces hexagonales) ou $\cos4\psi$ (faces carrées) \tierX, alors
qu'un $H$ axial y est constant \tierC.

\paragraph{Résultat, avec sa portée exacte.} Pour les classes d'essai radiale, axiale et « hiérarchique intrinsèque »,
un argument de projection \tierC, évalué en virgule flottante, force un déficit : la constante certifiable est au plus
$W_{\rm bcc}-\delta$ avec $\delta\ge8{,}751\times10^{-5}$ (hiérarchique) ou $\delta\ge1{,}058\times10^{-4}$ (axiale,
$1{,}18$ fois l'écart fcc--bcc $8{,}969\times10^{-5}$) ; ces valeurs sont \tierX. Pour des classes plus riches (qui lisent
les données des faces voisines), le négatif ne repose que sur une heuristique de comptage (environ 12 conditions pour 3
multiplicateurs, contre 3 pour 3 en dimension deux) \tierC{} et \textbf{n'est pas un théorème}. Rien n'est démontré sur
bcc.

% ---------------------------------------------------------------------
\section{Exercices}\label{ch2:sec:ex}

Les exercices marqués \textbf{(Lean)} renvoient à une solution de référence vérifiée par le noyau dans le dépôt
\tierA{} ; les solutions esquissées ici sont \tierC.

\begin{exercise}[témoin dipolaire]\label{ch2:ex:hd}
Soit $d>0$ et $h_d(t)=t^{-1/2}-2(t+d^2)^{-1/2}$. Calculer $h_d(d^2)$ et montrer qu'il est strictement négatif. En
déduire que $h_d$ n'est pas complètement monotone. Interpréter physiquement : quelle erreur de comptage de paires
produirait $h_d$ ?
\end{exercise}

\begin{exercise}[un mélange signé]\label{ch2:ex:signedmix}
Soit $\sigma(t)=e^{-t}-2e^{-2t}$. Montrer que $\sigma(\ln\frac32)=-\frac29$. Pourquoi le théorème~\ref{ch2:thm:laplace}
exige-t-il $m\ge0$ ? Exhiber une mesure signée dont la transformée de Laplace est $\sigma$.
\end{exercise}

\begin{exercise}[la borne de queue]\label{ch2:ex:diffle}
Redémontrer le lemme~\ref{ch2:lem:diffle} par les accroissements finis appliqués à $x\mapsto x^{-1/2}$ sur $[t,t+c]$.
Quelle constante obtient-on ? Comparer avec $\frac c2$.
\end{exercise}

\begin{exercise}[le noyau de la queue de Yukawa]
(a) Montrer que $e^x\ge x^2/2$ pour $x\ge0$. (b) En déduire $g_\kappa(t)\le\frac2{\kappa^2}t^{-3/2}$. (c) Montrer que,
dans le réseau $\Lat$, le nombre de points de la couronne $k\le|a|<k+1$ est $O(k)$ (indication : cellules de Voronoï
d'aire $1$, disjointes, contenues dans une couronne élargie), et conclure que $\sum_{a\ne0}|a|^{-3}<\infty$.
(d) Pourquoi l'argument échoue-t-il pour $\kappa=0$ ?
\end{exercise}

\begin{exercise}[moments du mélange de Keldysh]
Montrer que pour $u\ge0$, $c>0$ et $N\in\N$, $u^Ne^{-cu}\le N!/c^N$. En déduire que
$a\mapsto(4a)^{-n}e^{-c/(4a)}\,\frac{w(a)}{2a}$ est bornée sur $]0,1]$. (Lean : \lean{pow_mul_exp_neg_le}.)
\end{exercise}

\begin{exercise}[le piège $\top\le\top$]
Montrer que pour un réseau plan, $\sum_{0<|a|\le R}|a|^{-1}$ croît comme $R$. En déduire pourquoi l'on
\emph{s'attend} à ce que $\mathrm{latticeEnergy}(\mathcal K_{r_0})=+\infty$. Pourquoi cela ne rend-il pas le théorème
\lean{triangular_keldysh_optimal} faux, mais seulement vide ?
\end{exercise}

\begin{exercise}[quadrupôles]\label{ch2:ex:quad}
Soit $f$ CM et $c\ge0$, $g(t)=f(t)-f(t+c)$ et $\Delta_ah(t)=h(t)-h(t+a)$. (a) Vérifier que
$3f(t)-4f(t+c)+f(t+4c)=\Delta_cg(t)+\Delta_{2c}g(t)+\Delta_{3c}g(t)$. (b) En déduire, avec le point (v) de la
section~\ref{ch2:sec:lib}, que l'interaction de deux quadrupôles est CM. (c) Montrer que le champ d'un seul quadrupôle,
$t^{-1/2}-2(t+1)^{-1/2}+(t+4)^{-1/2}$, est négatif en $t=3$, en utilisant $\sqrt3>1{,}7$ et $\sqrt7>2{,}6$.
\end{exercise}

\begin{exercise}[CM ou Bernstein ?]
Montrer que : (a) $1/t$ est CM et positive mais n'est pas de Bernstein ; (b) $e^{-1/t}$ n'est pas CM (indication :
le signe de sa dérivée) ; (c) $t^\alpha$ est de Bernstein pour $0\le\alpha\le1$ mais pas pour $\alpha=2$. Expliquer
pourquoi (b) montre que, dans le théorème~\ref{ch2:thm:subord}, la fonction intérieure ne peut pas être n'importe
quelle fonction positive.
\end{exercise}

\begin{exercise}[l'hypothèse $\alpha+\beta>0$]\label{ch2:ex:signed}
Avec $G\equiv1$, $K=-1$, $d_0=d=1$ : vérifier les quatre hypothèses de \lean{signed_endpoints}, puis montrer que pour
$\alpha=\beta=-\frac\pi4$ on a $g_1(\tan\alpha)+g_1(\tan\beta)=-2<K$. Quelle étape de la démonstration du
théorème~\ref{ch2:thm:signed} utilise $\alpha+\beta>0$ ?
\end{exercise}

\begin{exercise}[arithmétique d'intervalles]
Avec $Q=100$, on code $a=[-1{,}25;0{,}33]$ par $(-125,33)$ et $b=[0{,}50;0{,}67]$ par $(50,67)$. (a) Calculer les quatre
produits d'extrémités et l'intervalle produit arrondi vers l'extérieur. (b) Avec la division entière euclidienne
$\lfloor\cdot\rfloor$, montrer que $x\mapsto-\lfloor -x/Q\rfloor$ calcule $\lceil x/Q\rceil$. (c) Pourquoi un arrondi
au plus proche serait-il incorrect pour un certificat ?
\end{exercise}

\begin{exercise}[\textbf{(Lean)} un témoin négatif]
En partant de \lean{h_neg_at : h d (d ^ 2) < 0} (pour $d>0$), démontrer en Lean \lstinline[style=chdeux]!¬ AdmissiblePotential (h d)!.
Indication : \lean{intro H}, puis utiliser la deuxième composante \lean{H.2.1} de la définition. Solution de référence :
\lean{formal_cert/Autoresearch/DipolarExciton.lean:113-116} \tierA.
\end{exercise}

\begin{exercise}[\textbf{(Lean)} corollaires de la bibliothèque]
(a) Énoncer en Lean que $t\mapsto\mathrm{riesz}\,s\,t-\mathrm{riesz}\,s\,(t+c)$ est CM pour $s,c\ge0$ et le démontrer en
une ligne à l'aide de \lean{riesz_cm} et \lean{sub_comp_add_const}. Solution de référence :
\lean{Corollaries.lean:33-35} \tierA. (b) Même question pour $e^{-at^\alpha}$, $a\ge0$, $0\le\alpha\le1$, à l'aide de
\lean{exp_neg_comp} et \lean{IsBernsteinOn.rpow}. Solution de référence : \lean{Corollaries.lean:59-61} \tierA.
(c) Pourquoi la condition \lstinline[style=chdeux]!(fun _ ht => Real.rpow_pos_of_pos ht α)! est-elle demandée ?
\end{exercise}

\label{ch2:fin}
